\chapter{Wyższe grupy homotopii i związek z homologią}
\label{ch:wyzsza-homotopia}
\index{grupa homotopii}

Grupa podstawowa z Rozdziału~\ref{ch:topologia-algebraiczna-podstawy}
rejestruje pętle. Wyższe grupy homotopii mierzą, czy można
ściągnąć do punktu odwzorowania całej sfery $S^n$. Homologia z
Rozdziału~\ref{ch:homologia-przestrzeni} bierze z kolei
zorientowane cykle i zapomina część informacji o sposobie ich
odwzorowania. Twierdzenie Hurewicza wskazuje dokładnie sytuację,
w której te dwa opisy mają ten sam pierwszy niezerowy wynik.
Twierdzenie Whiteheada mówi, kiedy wszystkie grupy homotopii
rozstrzygają typ homotopijny kompleksu CW.

W tym rozdziale pracujemy z przestrzeniami z wybranym punktem bazowym,
przede wszystkim z kompleksami CW. W argumentach komórkowych
wybieramy punkt bazowy będący wierzchołkiem. Gdy piszemy
$\pi_n(X)$ bez punktu bazowego, zakładamy spójność łukową i
wybieramy punkt bazowy; jego zmiana wzdłuż drogi daje izomorfizm, ale
sam wybór drogi może mieć znaczenie.

\section{Od pętli do odwzorowań sfer}
\label{sec:grupy-pi-n}

\begin{definicja}[Grupa $\pi_n$ i punkt bazowy]
\label{def:pi-n}
Niech $I=[0,1]$, $I^n=I\times\cdots\times I$, a
$\partial I^n$ niech oznacza sumę wszystkich jego ścian.
Dla przestrzeni z wybranym punktem $x_0$ kładziemy
\[
 \pi_n(X,x_0)
 =\big[(I^n,\partial I^n),(X,x_0)\big]_*
 =\big[(S^n,*),(X,x_0)\big]_*
 \qquad(n\geq1).
\]
Nawias oznacza klasy homotopii, w których
\emph{przez cały czas} brzeg kostki lub wyróżniony
punkt sfery jest wysyłany do $x_0$.
Utożsamienie pochodzi z $I^n/\partial I^n\cong S^n$.
Wyjaśnijmy przejście od kostki do dysku. Po przesunięciu
i przeskalowaniu $y=2s-(1,\ldots,1)$ kostka staje się
$[-1,1]^n$. Jej promień w normie maksimum to
$r=\|y\|_\infty$. Zachowując kierunek euklidesowy, wysyłamy
$y\ne0$ na $r\,y/\|y\|_2\in D^n$, a $0$ na $0$.
Norma obrazu wynosi $r$, więc brzeg przechodzi na $S^{n-1}$;
odwrotność dla $x\ne0$ to $\|x\|_2x/\|x\|_\infty$.
Oba wzory są ciągłe także w zerze, ponieważ normy obrazów
dążą do zera. Otrzymujemy homeomorfizm par
$(I^n,\partial I^n)\cong(D^n,S^{n-1})$.
Pozostaje $D^n/S^{n-1}\cong S^n$, wyprowadzone jawnie
w Uwadze~\ref{rem:iloraz-dysku-sfera}.
Dla $n=1$ otrzymujemy $\pi_1(X,x_0)$ z
Definicji~\ref{def:droga-petla}. Symbol $\pi_0(X)$
oznacza jedynie zbiór składowych łukowych,
zwykle bez działania grupowego.
\end{definicja}

Obraz odwzorowania $S^2\to X$ można sobie wyobrażać jako
elastyczną powłokę rozpiętą w przestrzeni; jej
ściągalność do punktu wymaga ciągłej deformacji
\emph{całego odwzorowania}, nie tylko zniknięcia klasy
homologii reprezentowanej przez to odwzorowanie. Sfera $S^n$ jest brzegiem dysku $D^{n+1}$,
a także jest homeomorficzna z brzegiem kostki $I^{n+1}$.

\begin{twierdzenie}[Działanie i przemienność]
\label{thm:pi-n-abelowe}
Dla $n\geq1$ klasy w $\pi_n(X,x_0)$ tworzą grupę.
Dla $n\geq2$ grupa ta jest abelowa.
\end{twierdzenie}
\begin{proof}
Jeśli $f,g:(I^n,\partial I^n)\to(X,x_0)$,
przedstawmy ich sumę sklejeniem w pierwszej
współrzędnej:
\[
 (f+_1g)(s_1,\ldots,s_n)=
 \begin{cases}
 f(2s_1,s_2,\ldots,s_n),&0\leq s_1\leq\tfrac12,\\
 g(2s_1-1,s_2,\ldots,s_n),&\tfrac12\leq s_1\leq1.
 \end{cases}
\]
Na płaszczyźnie sklejenia oba wzory dają
$x_0$; tak samo jest na zewnętrznym brzegu.
Homotopie $f$ lub $g$ zachowujące brzeg
sklejają się tym samym wzorem, więc działanie
nie zależy od reprezentantów. Łączność
wynika z liniowej zmiany podziału pierwszej
współrzędnej na trzy odcinki, dokładnie jak
w dowodzie Twierdzenia~\ref{thm:grupa-podstawowa}.
Stałe odwzorowanie jest elementem neutralnym:
reparametryzacje $r_-,r_+$ z tego samego dowodu stosujemy
do $s_1$, pozostawiając pozostałe współrzędne bez zmiany.
Odwrotność reprezentuje $\bar f(s_1,s')=f(1-s_1,s')$,
gdzie $s'=(s_2,\ldots,s_n)$. Jawna homotopia od
$f+_1\bar f$ do odwzorowania stałego to
\[
 K(s_1,s',t)=
 \begin{cases}
 f(2(1-t)s_1,s'),&s_1\leq\tfrac12,\\
 f(2(1-t)(1-s_1),s'),&s_1\geq\tfrac12.
 \end{cases}
\]
Pierwszy argument homotopii jest parametrem drogi
przebytej przez $f$ tam i z powrotem: dla $t=0$ rośnie
od $0$ do $1$ i maleje do $0$. Mnożnik $1-t$ skraca
osiągany przedział z $[0,1]$ do $[0,1-t]$.
Na styku obie gałęzie mają wartość $f(1-t,s')$.
Na brzegu kostki wartość jest zawsze $x_0$, a przy $t=1$
całe odwzorowanie ma wartość $f(0,s')=x_0$.
Zastąpienie $f$ przez $\bar f$ dowodzi też, że $\bar f+_1f$ jest zerem.

Gdy $n\geq2$, możemy niezależnie zdefiniować
drugie działanie $+_2$, sklejając w kierunku
$s_2$. Ma ono ten sam element neutralny
i tę samą konstrukcję homotopii. Podział
kwadratu $(s_1,s_2)$ na cztery ćwiartki
daje \emph{równość odwzorowań}, a nie tylko klas:
\[
 (a+_1 b)+_2(c+_1d)
   =(a+_2c)+_1(b+_2d).
\]
Teraz przechodzimy do klas homotopii i oznaczamy klasę stałą przez $0$.
Korzystając z neutralności $0$, otrzymujemy
\[
 a+_1b=(a+_2 0)+_1(0+_2 b)
      =(a+_1 0)+_2(0+_1 b)=a+_2b.
\]
Z kolei
\[
 a+_1b=(0+_2a)+_1(b+_2 0)
      =(0+_1b)+_2(a+_1 0)=b+_2a=b+_1a.
\]
Te ostatnie równości dotyczą klas; samo usuwanie stałych połówek
zmienia parametryzację odwzorowania.
To argument Eckmanna--Hiltona: dwa
zgodne sposoby sklejania wymuszają
przemienność. Użycie drugiej
współrzędnej jest niemożliwe dla
$n=1$, dlatego $\pi_1$ może być
nieabelowa.
\end{proof}

\begin{figure}[!htbp]
\centering
\begin{tikzpicture}[>=Stealth,line cap=round,line join=round]
 \fill[manifoldblue!8] (0,0) rectangle (4,4);
 \draw[manifoldblue!75!black,line width=1.4pt] (0,0) rectangle (4,4);
 \draw[manifoldblue!55] (2,0)--(2,4) (0,2)--(4,2);
 \fill[manifoldred!20] (.28,2.28) rectangle (1.72,3.72);
 \fill[manifoldgreen!22] (2.28,.28) rectangle (3.72,1.72);
 \node[font=\large,manifoldred!80!black] at (1,3) {$f$};
 \node[font=\large,manifoldgreen!70!black] at (3,1) {$g$};
 \node at (3,3) {$*$}; \node at (1,1) {$*$};
 \draw[->,manifoldgold!85!black,thick] (0,-.35)--(4,-.35)
  node[midway,below=4pt] {sklejanie w kierunku $s_1$};
 \draw[->,manifoldgold!85!black,thick] (-.35,0)--(-.35,4)
  ;
 \node[rotate=90,manifoldgold!85!black] at (-.8,2)
  {sklejanie w kierunku $s_2$};
 \node[font=\small,align=left,anchor=west,text width=4.2cm]
  at (4.6,2)
  {Puste pola są odwzorowaniem stałym. Oba kierunki
   sklejania opisują to samo odwzorowanie czterech
   ćwiartek; zamiana kolejności daje
   $[f]+[g]=[g]+[f]$.};
\end{tikzpicture}
\caption{Dwa niezależne kierunki sklejania w $\pi_2$.
W $\pi_n$ dla $n\geq2$ pozostałe współrzędne
traktujemy jako parametry.}
\label{fig:pi2-sklejanie}
\end{figure}

\begin{stwierdzenie}[Odwzorowania, iloczyny i homotopie]
\label{prop:pi-n-funktor}
Ciągłe odwzorowanie oparte $F:(X,x_0)\to(Y,y_0)$ wyznacza
homomorfizm $F_*:\pi_n(X,x_0)\to\pi_n(Y,y_0)$
przez $[f]\mapsto[F\circ f]$.
Odwzorowania homotopijne \emph{przez odwzorowania oparte}
indukują ten sam homomorfizm. Dla
$n\geq1$ i przestrzeni opartych $X,Y$:
\[
 \pi_n(X\times Y,(x_0,y_0))
 \cong \pi_n(X,x_0)\times\pi_n(Y,y_0).
\]
\end{stwierdzenie}
\begin{proof}
Złożenie $F\circ f$ nadal wysyła
$\partial I^n$ do $y_0$, a homotopia
$f_t$ daje homotopię $F\circ f_t$.
Ponieważ $F$ działa po sklejaniu
obu połówek dokładnie tak samo
jak przed nim, $F_*([f]+[g])
=F_*[f]+F_*[g]$. Homotopia
oparta $F_t$ daje homotopię
$F_t\circ f$ z nieruchomym
obrazem brzegu. Wreszcie odwzorowanie
$I^n\to X\times Y$ to dokładnie
para odwzorowań do $X$ i $Y$; to
samo dotyczy homotopii.
Projekcje i składanie pary
stanowią wzajemne odwrotności
podanego izomorfizmu.
\end{proof}

\begin{uwaga}[Zmiana punktu bazowego]
Niech $\gamma$ prowadzi od $x_0$ do $x_1$, a
$f:(I^n,\partial I^n)\to(X,x_1)$.
Umieścimy $f$ w środkowej kostce $[1/4,3/4]^n$,
a otaczający ją pas odwzorujemy wzdłuż $\gamma$.
Dla $s=(s_1,\ldots,s_n)$ połóżmy
$\rho(s)=\max_i|2s_i-1|$.
Na środkowej kostce każdą współrzędną musimy przenieść
z $[1/4,3/4]$ na $[0,1]$; wzór afiniczny to
$(s_i-1/4)/(3/4-1/4)=2s_i-1/2$.
Na pasie parametr drogi ma wynosić $1$ przy $\rho=1/2$
i $0$ przy $\rho=1$. Dla $\tau(\rho)=a+b\rho$
warunki $a+b/2=1$, $a+b=0$ dają $b=-2$, $a=2$.
Stąd definiujemy
\[
 (T_\gamma f)(s)=
 \begin{cases}
 f(2s_1-\tfrac12,\ldots,2s_n-\tfrac12),&\rho(s)\leq\tfrac12,\\
 \gamma(2-2\rho(s)),&\rho(s)\geq\tfrac12.
 \end{cases}
\]
Na styku obu części otrzymujemy $x_1=\gamma(1)$,
a na zewnętrznym brzegu $x_0=\gamma(0)$; lemat o sklejaniu
daje ciągłość. Ten sam wzór zastosowany do homotopii $f$
pokazuje niezależność od reprezentanta.
Homotopię drogi względem końców stosujemy tylko na pasie,
więc wynik zależy jedynie od klasy $\gamma$.

Składając dwie kostki z pasami, możemy przesunąć podział
tak, aby otrzymać pas wokół obu kostek; dwa przeciwne przejścia
drogi na wspólnej ścianie kasują się przez homotopię
drogi przebytej tam i z powrotem, jak w dowodzie
Twierdzenia~\ref{thm:grupa-podstawowa}.
Dlatego $T_\gamma$ zachowuje działanie na klasach.
Dla $n=1$ ten wzór przedstawia wprost pętlę
$\gamma*f*\bar\gamma$, z nieistotnymi zmianami parametryzacji.
Zagnieżdżenie dwóch pasów odpowiada złożeniu dróg:
jeśli $\eta$ prowadzi od $x_1$ do $x_2$, to
$T_\gamma\circ T_\eta=T_{\gamma*\eta}$ na klasach homotopii.
Pas dla drogi stałej można zwęzić do brzegu, więc indukuje identyczność.
Ponieważ $\gamma*\bar\gamma$ i $\bar\gamma*\gamma$ są homotopijne
do dróg stałych względem końców, $T_{\bar\gamma}$ jest odwrotnością $T_\gamma$.
Otrzymujemy zatem izomorfizm
$T_\gamma:\pi_n(X,x_1)\to\pi_n(X,x_0)$.
Dwie różne drogi mogą dać różne izomorfizmy; dla pętli
otrzymujemy \emph{działanie} $\pi_1(X,x_0)$ na $\pi_n(X,x_0)$.
W przypadku przestrzeni jednospójnej wybór drogi
nie powoduje tej niejednoznaczności.
Jeśli $F_t:X\to Y$ jest homotopią bez warunku oparcia i
$\gamma(t)=F_t(x_0)$, to
$T_\gamma\circ(F_1)_*=(F_0)_*$.
Istotnie, w powyższym wzorze używamy kolejno drogi
$\gamma_t(u)=\gamma(tu)$ i odwzorowania $F_t\circ f$.
Daje to homotopię opartą w $F_0(x_0)$; przy $t=0$ pas jest stały,
a przy $t=1$ otrzymujemy $T_\gamma(F_1\circ f)$.
W szczególności zwykła równoważność homotopijna indukuje
izomorfizmy wszystkich grup homotopii, z odpowiednio dobranymi bazami.
Dowód Wniosku~\ref{cor:pi1-homotopia} używał właśnie
funktorialności i zmiany bazy, więc ten sam argument działa dla każdego $n$.
\end{uwaga}

\section{Grupy względne i długi ciąg homotopii}
\label{sec:pi-wzgledne}

Tak jak $H_n(X,A)$ mierzy cykle,
których brzeg może trafić do $A$,
tak względna homotopia mierzy
dyski w $X$ z brzegiem w $A$.
Brzeg dysku może się przy tym
przesuwać w $A$ podczas homotopii.

\begin{definicja}[Względne grupy homotopii]
\index{wzgledne grupy homotopii@Względne grupy homotopii}
\label{def:pi-wzgledne}
Niech $x_0\in A\subseteq X$.
W kostce $I^n$ wyróżnijmy
dolną ścianę $B=I^{n-1}\times\{0\}$,
a przez $J$ oznaczmy sumę
pozostałych ścian kostki $I^n$.
Klasa względna jest klasą odwzorowań
\[
 F:(I^n,\partial I^n,J)
       \longrightarrow(X,A,x_0);
\]
zatem $F(\partial I^n)\subseteq A$
i $F(J)=x_0$. Homotopie spełniają
te same warunki. Ich zbiór oznaczamy
$\pi_n(X,A,x_0)$.
W równoważnym modelu są to
oparte odwzorowania trójek
$(D^n,S^{n-1},*)\to(X,A,x_0)$.
Dla $n\geq2$ składanie w pierwszej
współrzędnej daje grupę; dla
$n\geq3$ jest ona abelowa.
Dla $n=1$ pozostaje \emph{zbiór
z wyróżnionym elementem},
ponieważ nie ma wolnej współrzędnej
do składania.
\end{definicja}

Przy $n\geq2$ płaszczyzna sklejania $s_1=1/2$ jest obrazem
ścian $s_1=0,1$, należących do $J$, więc oba wzory mają na niej
wartość $x_0$. Homotopie reparametryzacji i kasowania odwrotności
z Twierdzenia~\ref{thm:pi-n-abelowe} zachowują teraz dolną ścianę w $A$,
a pozostałe ściany w $x_0$; dają więc aksjomaty grupy względnej.
Jeśli $n\geq3$, współrzędna $s_2$ również nie jest ostatnią współrzędną,
zatem drugi sposób sklejania zachowuje te warunki i dowodzi przemienności.

W przeciwieństwie do $\pi_n(X)$,
względna $\pi_2(X,A)$ może być
nieprzemienna. Drugi kierunek
sklejania zajmuje wówczas rolę
wyróżnionej ściany, a argument
z Rysunku~\ref{fig:pi2-sklejanie}
nie ma zastosowania.

\begin{figure}[!htbp]
\centering
\begin{tikzpicture}[>=Stealth,line cap=round,line join=round]
 \filldraw[rounded corners=8pt,fill=manifoldblue!7,draw=manifoldblue!60!black]
  (0,0) rectangle (8,3.4);
 \fill[manifoldblue!18] (2.5,1.65) ellipse (1.65 and 1.05);
 \draw[manifoldgreen!28,line width=12pt] (2.5,1.65) ellipse (1.65 and 1.05);
 \draw[manifoldred,line width=1.4pt] (2.5,1.65) ellipse (1.65 and 1.05);
 \fill[manifoldgold!80!black] (.85,1.65) circle(2.6pt);
 \node[left,font=\small] at (.78,1.65) {$x_0$};
 \node[manifoldblue!75!black] at (2.5,1.65) {$F(\operatorname{int}D^2)$};
 \node[anchor=west,manifoldgreen!60!black] at (5.1,2.5) {$A$ (zielony pas)};
 \draw[->,manifoldgreen!60!black] (5,2.5)--(3.85,2.32);
 \node[anchor=west,manifoldred!80!black] at (5.1,1.05) {$F(S^1)\subset A$};
 \draw[->,manifoldred!80!black] (5,1.05)--(4.05,1.3);
 \node[manifoldblue!75!black] at (7.65,3.05) {$X$};
\end{tikzpicture}
\caption{Przykład odwzorowania względnego $(D^2,S^1,*)\to(X,A,x_0)$.
Czerwony obraz brzegu leży w zielonym pasie $A$, podczas gdy
obraz wnętrza dysku może wyjść poza $A$. W ogólnej definicji
nie wymaga się różnowartościowości $F$.}
\label{fig:pi-wzgledny-dysk}
\end{figure}

\begin{twierdzenie}[Długi ciąg homotopii pary]
\label{thm:les-pi}
Dla $x_0\in A\subseteq X$ otrzymujemy
naturalny ciąg
\begin{equation}\label{eq:les-pi}
 \cdots\to\pi_n(A,x_0)\xrightarrow{i_*}
 \pi_n(X,x_0)\xrightarrow{j_*}
 \pi_n(X,A,x_0)\xrightarrow{\partial}
 \pi_{n-1}(A,x_0)\xrightarrow{i_*}
 \pi_{n-1}(X,x_0)\to\cdots .
\end{equation}
Jest on dokładny: obraz każdej
strzałki to jądro następnej.
Odwzorowanie $\partial$ ogranicza
reprezentujący dysk do jego
wyróżnionej ściany $B$.
Na końcu ciągu występują
zbiory $\pi_1(X,A,x_0)$ i
$\pi_0(A),\pi_0(X)$; dokładność
oznacza tam przeciwobraz
wyróżnionego elementu.
\end{twierdzenie}
\begin{proof}
W modelu kostkowym $i_*$ jest
odwzorowaniem indukowanym przez inkluzję $A\to X$, a $j_*$
traktuje odwzorowanie z całym brzegiem
w $x_0$ jako odwzorowanie względne.
Wzór $\partial[F]=[F|_B]$
jest poprawny, ponieważ
$\partial B\subset J$ i
$F(J)=x_0$. Homotopia
względna ogranicza się
do homotopii opartych
odwzorowań na $B$, więc
$\partial$ nie zależy
od reprezentanta.

Sprawdźmy trzy kolejne miejsca.
\emph{W $\pi_n(X)$.} Jeśli $f$ ma cały obraz w $A$, homotopia
$f_t(s',r)=f(s',(1-t)r+t)$ jest względną homotopią do stałej;
tutaj $s'\in I^{n-1}$, a $r$ jest ostatnią współrzędną.
Odwrotnie, niech $H(s',r,t)$ będzie względną homotopią
od $f$ do stałej. Ściana $g(s',t)=H(s',0,t)$ leży w $A$
i ma cały brzeg w $x_0$, więc reprezentuje element $\pi_n(A)$.
W kwadracie $(r,t)$ porównujemy drogi od $(0,0)$ do $(1,1)$:
dolny, a potem prawy bok oraz lewy, a potem górny bok.
Są homotopijne względem końców przez liniową interpolację w kwadracie.
Po złożeniu z $H$ dają odpowiednio $f$ i $g$ ze stałym dodatkowym
odcinkiem; na brzegu współrzędnych $s'$ wszystko pozostaje w $x_0$.
Usunięcie stałego odcinka jest reparametryzacją ostatniej współrzędnej,
więc nie porusza brzegu kostki. Stąd $[f]=i_*[g]$.

\emph{W $\pi_n(X,A)$.} Dysk
pochodzący z $\pi_n(X)$
ma stały brzeg, więc jego
$\partial$ jest zerem.
Odwrotnie, jeśli $[F|_B]=0$ w $\pi_{n-1}(A)$ dla $n\geq2$,
wybieramy opartą nullhomotopię
$h:I^{n-1}\times I\to A$ z $h(s',0)=F(s',0)$ i $h(s',1)=x_0$.
Trzeba włączyć tę deformację brzegu do całego dysku.
Dla $\tau\in I$ doklejamy na dole pas grubości $\tau/2$.
Na pasie czas homotopii $h$ ma maleć od $\tau$ przy $r=0$
do $0$ przy $r=\tau/2$, zatem wynosi $\tau-2r$.
Pozostały odcinek $[\tau/2,1]$ przeskalowujemy na $[0,1]$:
odejmujemy jego początek $\tau/2$ i dzielimy przez długość $1-\tau/2$.
Otrzymujemy
\[
 K_\tau(s',r)=
 \begin{cases}
 h(s',\tau-2r),&0\leq r\leq\tau/2,\\
 F\left(s',\dfrac{r-\tau/2}{1-\tau/2}\right),&\tau/2\leq r\leq1.
 \end{cases}
\]
Na styku obie wartości są równe $F(s',0)$, a mianownik
jest co najmniej $1/2$, więc wzór daje ciągłą homotopię,
także dla $\tau=0$. Dolna ściana pozostaje w $A$,
a $J$ w $x_0$. Mamy $K_0=F$, natomiast $K_1$ ma już
\emph{cały} brzeg w $x_0$, ponieważ na dole przyjmuje $h(s',1)=x_0$.
Jest więc reprezentantem klasy w $\pi_n(X)$,
której obraz przez $j_*$ jest równy $[F]$.
Geometrycznie jest to sklejenie dysku $F$ z wypełnieniem jego brzegu w $A$;
wzór na $K_\tau$ uzasadnia równość otrzymanych klas względnych.

\emph{W $\pi_{n-1}(A)$.} Brzeg
dysku względnego wypełnia
się w $X$, więc ginie
po inkluzji do $X$.
Jeśli oparte odwzorowanie $g:I^{n-1}\to A$ jest
zerowe w $\pi_{n-1}(X)$, wybieramy opartą homotopię
$G(s',r)$ od $g$ do stałej. Jest ona zarazem odwzorowaniem względnym:
na dole ma wartość $g(s')\in A$, na górze $x_0$, a na pozostałych
ścianach $x_0$, ponieważ homotopia jest oparta.
Zatem $\partial[G]=[g]$. W drugą stronę ten sam wzór pokazuje,
że brzeg każdego dysku względnego ginie po inkluzji do $X$.
Te trzy rachunki powtarzają się w sąsiednich stopniach.
Na końcu ciągu element $\pi_1(X,A,x_0)$ jest drogą
od pewnego $a\in A$ do $x_0$, a $\partial$ przypisuje mu
składową $a$ w $A$. Trafia ona do wyróżnionej składowej $X$.
Odwrotnie, punkt tej składowej można połączyć drogą z $x_0$.
Jeśli $a$ leży w składowej $x_0$ już w $A$, dopinamy drogę
z $x_0$ do $a$ w $A$ i otrzymujemy pętlę w $X$ reprezentującą
daną klasę względną. To dowodzi dokładności także dla zbiorów wskazanych.
Wszystkie konstrukcje zachowują się przy składaniu z odwzorowaniem par,
co daje naturalność ciągu.
\end{proof}

\begin{definicja}[Spójność homotopijna]
\index{spojnosc homotopijna@Spójność homotopijna}
\label{def:spojnosc-homotopijna}
Przestrzeń $X$ jest
\emph{$r$-spójna} ($r\geq0$),
jeśli jest łukowo spójna
i $\pi_i(X)=0$ dla $1\leq i\leq r$.
W szczególności $1$-spójna
oznacza jednospójną.
Para $(X,A)$ jest
\emph{$r$-spójna}, jeśli
każda składowa łukowa $X$
spotyka $A$ oraz
$\pi_i(X,A)=0$ dla
$1\leq i\leq r$ przy każdym wyborze punktu bazowego w $A$
(w stopniu $1$ jest to
trywialność zbioru
z wyróżnioną klasą).
\end{definicja}

Dokładność~\eqref{eq:les-pi}
pokazuje, że dla $r$-spójnej
pary inkluzja $A\hookrightarrow X$
indukuje izomorfizm na
$\pi_i$ dla $i<r$ oraz
surjekcję na $\pi_r$.
Rzeczywiście, w fragmencie
$\pi_{i+1}(X,A)\to\pi_i(A)\to\pi_i(X)\to\pi_i(X,A)$
oba zewnętrzne wyrazy są zerowe dla $1\leq i<r$;
przy $i=r\geq1$ znika prawy wyraz, co daje surjektywność.
Na poziomie składowych surjektywność wynika z tego, że każda
składowa $X$ spotyka $A$. Jeśli dodatkowo $r\geq1$, droga w $X$
łącząca dwa punkty $A$ reprezentuje zerową klasę względną,
więc jej początek można przesunąć w $A$ do jej końca;
stąd także bijekcja $\pi_0(A)\to\pi_0(X)$.
Wymiar komórek CW pozwoli
za chwilę uzyskać takie
zanikanie bez liczenia odwzorowań
pojedynczo.

\section{Homotopie na komórkach CW}
\label{sec:pi-n-cw}

W rozdziale o homologii komórki dawały
generatory grup łańcuchów. W homotopii
ich wymiary pozwalają \emph{usuwać
przeszkody do deformacji}. Żeby
deformować jedną komórkę bez psucia
tego, co zrobiliśmy na jej brzegu,
potrzebujemy przedłużania homotopii.

\begin{lemat}[Zwarte obrazy i homotopie w CW]
\label{lem:zwartosc-cw-pi}
Każdy zwarty podzbiór kompleksu CW leży w skończonym podkompleksie.
Ponadto odwzorowanie $H:X\times I\to Y$ jest ciągłe, jeśli jego złożenia
z każdym $\chi_\alpha\times\id_I:D^{d_\alpha}\times I\to X\times I$
są ciągłe, gdzie $\chi_\alpha$ są odwzorowaniami charakterystycznymi komórek.
\end{lemat}
\begin{proof}
Gdyby zwarty zbiór $C$ spotykał nieskończenie wiele otwartych komórek,
wybralibyśmy z każdej z nich jeden punkt $c_\alpha\in C$.
Każda domknięta komórka spotyka tylko skończenie wiele spośród tych punktów.
Każdy podzbiór $E$ zbioru wybranych punktów ma więc domknięty
przeciwobraz przez każde odwzorowanie charakterystyczne: jest to skończona
suma przeciwobrazów punktów w przestrzeni Hausdorffa.
Topologia słaba mówi, że $E$ jest domknięty w $X$.
Zatem zbiór wszystkich $c_\alpha$ jest domkniętym, dyskretnym
podzbiorem zwartego $C$, co jest niemożliwe, jeśli jest nieskończony.
Istotnie, jako domknięty podzbiór $C$ byłby zwarty, natomiast
pokrycie dyskretnego zbioru przez jego jednoelementowe podzbiory
nie ma skończonego podpokrycia.
Do skończenie wielu spotkanych komórek dołączamy ich ściany;
warunek skończoności domknięć i malejące wymiary dają skończony podkompleks.

Topologia słaba oznacza, że surjekcja
$q:\bigsqcup_\alpha D^{d_\alpha}\to X$ jest ilorazowa.
Iloczyn dowolnego odwzorowania ilorazowego $q:E\to X$ z $\id_I$ też jest ilorazowy.
Aby to sprawdzić, niech przeciwobraz $W\subseteq X\times I$ będzie otwarty.
Dla $e\in q^{-1}(x)$ przekrój
$\{s\in I:(e,s)\in(q\times\id_I)^{-1}(W)\}$ jest otwarty w $I$
i jest równy $\{s:(x,s)\in W\}$.
Możemy więc dla $(x,t)\in W$ wybrać zwarte otoczenie $J$ punktu $t$
w tym przekroju, zatem $\{x\}\times J\subset W$. Zbiór
$V=\{y:\{y\}\times J\subset W\}$ ma otwarty przeciwobraz przez $q$:
dla każdego $e$ nad $V$ skończone pokrycie zwartego $J$
prostokątami zawartymi w przeciwobrazie $W$ daje otoczenie $e$
o tej samej własności; bierzemy przecięcie skończenie wielu
otoczeń $e$ występujących w tych prostokątach.
Z ilorazowości $q$ wynika, że $V$ jest otwarty, a
$V\times\operatorname{int}_I J\subseteq W$ jest otoczeniem $(x,t)$.
To dowodzi tezy o iloczynie, a więc także kryterium ciągłości $H$.
\end{proof}

Z lematu wynika również rozszerzenie homologii komórkowej
z Rozdziału~\ref{ch:homologia-przestrzeni} na dowolne CW i pary CW.
Dla pary skończonych CW definiujemy najpierw
$C_\bullet^{\mathrm{CW}}(X,A)=C_\bullet^{\mathrm{CW}}(X)/C_\bullet^{\mathrm{CW}}(A)$:
jest to kompleks wolny na komórkach spoza $A$, z brzegiem modulo komórki $A$.
Krótki ciąg tych kompleksów daje długi ciąg homologii.
Naturalne porównanie ze zwykłym ciągiem pary i lemat pięciu
przenoszą izomorfizmy dla $A,X$ na grupy względne.
Każdy singularny łańcuch, jego brzeg oraz świadectwo homologiczności
są skończonymi sumami odwzorowań zwartych sympleksów, więc mieszczą się
w skończonym podkompleksie. Cykl pochodzi z takiego podkompleksu,
a relacja między cyklami jest już prawdziwa w pewnym większym,
również skończonym. Tak samo jest dla skończonych sum komórek
i dla par: przecięcie skończonego podkompleksu z podkompleksem $A$
jest skończonym podkompleksem. Izomorfizmy dla modeli skończonych
są zgodne z inkluzjami, bo pochodzą z naturalnych ciągów par.
Przechodzą więc na izomorfizmy dla całego CW.
Grupy łańcuchów komórkowych są przy tym \emph{sumami prostymi},
a nie iloczynami: każda kombinacja nadal ma skończony nośnik.

\begin{lemat}[Przedłużanie homotopii dla par CW]
\label{lem:hep-cw}
Niech $A\subseteq X$ będzie podkompleksem
CW. Jeśli $f:X\to Y$ i homotopia
$h:A\times I\to Y$ spełniają
$h(a,0)=f(a)$, to istnieje
$H:X\times I\to Y$ takie, że
$H(x,0)=f(x)$ oraz $H|_{A\times I}=h$.
W szczególności homotopię
zadaną na brzegu dołączonego
dysku można przedłużyć do
całego dysku.
\end{lemat}
\begin{proof}
Najpierw rozważmy jedną komórkę
$D^k$, której brzeg $S^{k-1}$
leży w już zbudowanej części.
Na
\[
 (D^k\times\{0\})\cup
 (S^{k-1}\times I)
\subset D^k\times I
\]
mamy dane: $f$ na dole
i homotopię $h$ na boku.
Ta podprzestrzeń jest
retraktem całego cylindra. Poprowadźmy promień z punktu
$(0,2)$, leżącego nad cylindrem, przez $(x,t)$.
Jego parametryzacja to
$(0,2)+\lambda((x,t)-(0,2))=(\lambda x,2-\lambda(2-t))$;
punkt wyjściowy $(x,t)$ odpowiada $\lambda=1$.
Zwiększamy $\lambda$, aż po raz pierwszy trafimy na bok
lub na dół. Pisząc $r=\|x\|$, dostajemy odpowiednio równania
$\lambda r=1$ oraz $2-\lambda(2-t)=0$.
Wybieramy mniejszy z dwóch parametrów trafienia, czyli
\[
 \lambda(x,t)=
 \min\left\{\frac1r,\frac2{2-t}\right\},
 \qquad
 R(x,t)=\bigl(\lambda x,\,
                  2-\lambda(2-t)\bigr),
\]
przy $r=0$ przyjmujemy
$1/r=+\infty$. Dla $r\leq1/2$ minimum wynosi $2/(2-t)$,
więc wzór jest ciągły także przy $r=0$. Wartość
$R(x,t)$ leży albo na boku
$\|\lambda x\|=1$, albo
na dole z drugą współrzędną
równą $0$. Ponadto $\lambda\geq1$ i
$\lambda\leq2/(2-t)$, więc druga współrzędna należy do $[0,t]\subseteq I$.
Jeśli $t=0$
lub $r=1$, mamy $\lambda=1$,
więc $R$ pozostawia dane
punkty nieruchome.
Złożenie danych z $R$
przedłuża homotopię
na $D^k\times I$.

Dołączajmy teraz komórki
$X\setminus A$ według
wymiarów. Na każdym kroku
stosujemy poprzednią
konstrukcję na wszystkich
dyskach danego wymiaru;
na ich brzegach odwzorowania
już się zgadzają. Ponieważ
topologia CW jest zadana
przez odwzorowania charakterystyczne,
sklejone homotopie są
ciągłe. Dla nieskończenie wielu komórek uzasadnieniem ciągłości
z parametrem czasu jest Lemat~\ref{lem:zwartosc-cw-pi},
zastosowany do wszystkich odwzorowań $D^k\times I\to Y$.
Dla $0$-komórek wybieramy homotopię stałą równą wartości $f$;
wzór na $R$ jest potrzebny dla $k\geq1$.
\end{proof}

\begin{twierdzenie}[Przybliżenie komórkowe]
\label{thm:przyblizenie-komorkowe}
Każde odwzorowanie $f:X\to Y$ kompleksów CW
jest homotopijne z odwzorowaniem
\emph{komórkowym}, czyli takim,
które spełnia
$f(X^k)\subseteq Y^k$ dla
każdego $k$. Jeśli na
podkompleksie $A\subseteq X$
odwzorowanie już jest komórkowe,
homotopię można wybrać
stałą na $A$.
\end{twierdzenie}
\begin{proof}
Wyjaśnijmy lokalny krok geometryczny. Niech
$u:D^k\to Y^m$ wysyła brzeg do $Y^{k-1}$, gdzie $m>k$.
Obraz jest zawarty w skończonym podkompleksie na mocy
Lematu~\ref{lem:zwartosc-cw-pi}. W jednej otwartej $m$-komórce
używamy współrzędnych $\operatorname{int}D^m$ i oznaczamy przez $v$
współrzędnościowy zapis $u$ na otwartym przeciwobrazie tej komórki.
Wybierzmy $0<a<b<c<1$. Zbiór $K=v^{-1}(\overline B_b)$
jest zwarty i leży we wnętrzu $D^k$: obraz $\overline B_b$
jest zwarty i domknięty w przestrzeni Hausdorffa $Y$.
Wybieramy skończone otoczenie wielościenne $P$ zbioru $K$
takie, że $K\subset\operatorname{int}P\subset P\subset v^{-1}(B_c)$.
Można je uzyskać jako sumę domkniętych kostek dostatecznie drobnej siatki
spotykających $K$. Triangulujemy $P$ i przez interpolację wartości $v$
w wierzchołkach dostajemy odcinkowo afiniczne $w$ z
$\|w-v\|<\varepsilon$, gdzie $\varepsilon<\min(b-a,1-c)$.
Istnienie takiego przybliżenia wynika z jednostajnej ciągłości $v$ na $P$:
na każdym małym sympleksie interpolujemy wartości różniące się o mniej niż $\varepsilon$.

Niech $\eta:D^k\to[0,1]$ będzie ciągła, równa $1$ na $K$
i $0$ poza $\operatorname{int}P$. Na $P$ stosujemy homotopię
$v_t=v+t\eta(w-v)$, a poza $P$ nic nie zmieniamy.
Na $P$ nierówność trójkąta daje
$\|v_t\|\leq\|v\|+t\eta\|w-v\|<c+\varepsilon<1$.
Wszystkie wartości pozostają więc we wnętrzu komórki,
a poza $P$ homotopia niczego nie zmienia, w szczególności na brzegu dysku.
W $P\setminus K$ mamy $\|v\|>b$ i
$\|v_t\|\geq\|v\|-t\eta\|w-v\|>b-\varepsilon>a$;
poza $P$ ten sam wniosek wynika z $v_t=v$.
W $K$ końcowe odwzorowanie jest równe $w$.
Jego obraz jest zawarty w skończonej sumie podprzestrzeni afinicznych
wymiaru co najwyżej $k<m$, która nie zawiera otwartej kuli $B_a$.
Istnieje więc $z\in B_a$ pominięty przez całe nowe odwzorowanie.

Promieniowa retrakcja $D^m\setminus\{z\}$ na $\partial D^m$
zostawia brzeg nieruchomy, więc po sklejeniu komórki daje homotopię
usuwającą jej wnętrze z obrazu $u$.
Wyznaczmy punkt trafienia: dla jednostkowego kierunku $\xi$
warunek $\|z+R\xi\|^2=1$ oznacza
$R^2+2\langle z,\xi\rangle R+\|z\|^2-1=0$.
Ponieważ $\|z\|<1$, pierwiastki mają przeciwne znaki.
Dodatni pierwiastek wynosi
\[
 R_z(\xi)=-\langle z,\xi\rangle+
 \sqrt{\langle z,\xi\rangle^2+1-\|z\|^2}.
\]
Jest ciągły, dodatni i nie większy niż $1+\|z\|$,
bo jest odległością między $z$ i punktem sfery jednostkowej.
Na promieniu $z+r\xi$, gdzie $0<r\leq R_z(\xi)$,
homotopia ma zatem wzór
\[
 z+r\xi\longmapsto z+\bigl((1-t)r+tR_z(\xi)\bigr)\xi.
\]
Interpolujemy liniowo od odległości $r$ do odległości trafienia.
Nie używamy wzoru w $z$, który został pominięty przez odwzorowanie.
Na brzegu $r=R_z(\xi)$ nic się nie zmienia,
więc homotopia jest zgodna z odwzorowaniem dołączającym
i z identycznością na pozostałych komórkach.
Odwzorowanie $u$ pomija $z$, zatem można je złożyć z tą homotopią.
Ciągłość po sklejeniu, również z parametrem czasu,
wynika z ilorazowego kryterium z Lematu~\ref{lem:zwartosc-cw-pi};
usunięcie zamkniętego punktu $z$ ogranicza iloraz do nasyconego zbioru otwartego.
Usuwamy kolejno komórki najwyższego wymiaru i schodzimy z wymiarem do $k$.
Pracujemy w skończonym podkompleksie zawierającym obraz i domknięcia
spotkanych komórek, zatem proces kończy się po skończenie wielu krokach.

Zaczynamy od $X^0$:
obraz każdego wierzchołka
przesuwamy do odpowiedniego
wierzchołka w jego składowej
$Y$, a na $A$ niczego
nie zmieniamy. Załóżmy
indukcyjnie, że
$f(X^{k-1})\subseteq Y^{k-1}$.
Dla każdej $k$-komórki
stosujemy lokalny krok
do jej odwzorowania
charakterystycznego:
usuwamy po kolei obrazy
komórek celu o wymiarach
$m>k$, zachowując
jej brzeg. Dostajemy
odwzorowanie tej komórki do
$Y^k$. Dzięki
Lematowi~\ref{lem:hep-cw}
homotopię na dotychczas
zbudowanej części
przedłużamy na całe $X$
i przechodzimy do
następnego wymiaru.
Podane kroki wykonuje
się względnie $A$.
Etapowi wymiaru $k$ przydzielamy przedział
$[1-2^{-k},1-2^{-k-1}]$, a w chwili $1$ bierzemy odwzorowanie końcowe.
Każdy ustalony szkielet pozostaje nieruchomy po zakończeniu swojego etapu.
Przy przedłużaniu kolejnego etapu zachowujemy bowiem homotopię stałą
na podkompleksie $A\cup X^{k-1}$.
Na każdym $D^k\times I$ mamy więc skończenie wiele sklejonych homotopii
i stały odcinek końcowy. Lemat~\ref{lem:zwartosc-cw-pi}
zapewnia ciągłość całej homotopii, także w chwili $1$.
\end{proof}

\begin{uwaga}[Przybliżenie w równych wymiarach]
\label{rem:przyblizenie-rowne-wymiary}
Ta sama lokalna konstrukcja dla $k=m$ pozwala wybrać $z\in B_a$
poza obrazami ścian triangulacji $P$ i tych $m$-sympleksów,
na których $w$ ma rząd mniejszy od $m$. Jest to skończona suma
zbiorów o wymiarze mniejszym od $m$. Każdy przeciwobraz $z$
leży wtedy we wnętrzu jednego sympleksu, gdzie $w$ jest afinicznym homeomorfizmem.
Przeciwobrazów jest skończenie wiele, a każdy ma znak wyznacznika $+1$ lub $-1$.
Nie leżą one na brzegu $K$: tam $\|v\|=b$, więc
$\|w\|>b-\varepsilon>a>\|z\|$.
W otoczeniu każdego przeciwobrazu mamy zatem $\eta=1$,
a końcowe odwzorowanie jest dokładnie afinicznym $w$.
Możemy wokół $z$ wybrać małą kulę, której przeciwobrazy są rozłącznymi
dyskami, odwzorowanymi homeomorficznie na tę kulę.
To wynika również ze zwartości: poza małymi otoczeniami znalezionych
przeciwobrazów obraz omija całe dostatecznie małe otoczenie $z$.
Ten wariant będzie użyty w dowodzie o pierwszych komórkach.
\end{uwaga}

\begin{wniosek}[Szkielet widziany przez $\pi_i$]
\label{cor:pi-szkielety}
Dla kompleksu CW $X$ para
$(X,X^n)$ jest $n$-spójna.
Inkluzja $X^n\hookrightarrow X$
daje więc izomorfizm na
$\pi_i$ dla $i<n$
i surjekcję dla $i=n$.
Jeśli $1\leq i<k$, to
$\pi_i(S^k)=0$.
\end{wniosek}
\begin{proof}
Reprezentant klasy
$\pi_i(X,X^n)$ jest
odwzorowaniem $(D^i,\partial D^i)\to
(X,X^n)$, gdzie $i\leq n$.
Najpierw przybliżamy komórkowo odwzorowanie brzegu do $X^n$,
zachowując punkt bazowy, i przedłużamy tę homotopię na dysk
Lematem~\ref{lem:hep-cw}. Dopiero teraz brzeg jest komórkowy,
więc względne przybliżenie na dysku przesuwa cały obraz do
$X^i\subseteq X^n$. Klasa względna znika: jeśli całe
$F(I^{i-1}\times I)$ leży w $X^n$, to
$F_t(s',r)=F(s',(1-t)r+t)$ jest dozwoloną homotopią
względną od $F$ do stałej. Na dolnej ścianie może zmieniać wartości
w $X^n$, a pozostałe ściany nadal mają wartość $x_0$.
Każda składowa $X$ zawiera wierzchołek: punkt komórki można
połączyć z jej brzegiem i powtarzać to w malejących wymiarach.
To daje także stwierdzenie o składowych i przypadek $n=0$.
Długi ciąg~\eqref{eq:les-pi}
daje twierdzenie o
inkluzji szkieletu.
Sfera $S^k$ ma strukturę CW
z jednym wierzchołkiem
i jedną $k$-komórką,
więc jej $i$-szkielet
jest punktem dla $0\leq i<k$.
Każde oparte odwzorowanie
$S^i\to S^k$ jest
homotopijne do odwzorowania
komórkowego i wtedy
jego obraz jest punktem.
\end{proof}

\begin{stwierdzenie}[Wyższe grupy i nakrycia]
\label{prop:pi-n-nakrycie}
Jeśli $p:\widetilde X\to X$
jest nakryciem, $\widetilde x_0$
leży nad $x_0$, to
\[
 p_*:\pi_n(\widetilde X,\widetilde x_0)
 \xrightarrow{\ \cong\ }
 \pi_n(X,x_0)\qquad(n\geq2).
\]
\end{stwierdzenie}
\begin{proof}
Sfera $S^n$ dla $n\geq2$
jest jednospójna na mocy
Wniosku~\ref{cor:pi-szkielety}
(dla $n=2$ znika $\pi_1$,
a dla większych wymiarów
również). Każde oparte
odwzorowanie $S^n\to X$ podnosi
się więc jednoznacznie
po wskazaniu obrazu
punktu bazowego,
zgodnie z kryterium
z Twierdzenia~\ref{thm:kryterium-podnoszenia}.
Daje to surjektywność $p_*$.
Jeśli dwa podniesienia
po zrzutowaniu są
homotopijne przez odwzorowania oparte,
ich homotopię
$S^n\times I\to X$ podnosimy z kryterium podnoszenia:
$S^n\times I$ jest łukowo i lokalnie łukowo spójna oraz jednospójna,
bo retrahuje się na $S^n$. Ustalamy obraz $(\ast,0)$.
Ścieżka punktu bazowego leży w dyskretnym włóknie, więc jest stała,
a ograniczenie do $S^n\times\{0\}$ jest zadanym podniesieniem.
Jednoznaczność podnoszenia
sprawia, że na końcu
otrzymujemy drugie
podniesienie, ponieważ
oba zgadzają się w
punkcie bazowym.
To daje iniektywność.
\end{proof}

\begin{przyklad}[Okrąg i torus nie mają wyższych $\pi_n$]
\label{prz:pi-torus}
Nakrycia $\R\to S^1$ i
$\R^2\to T^2=S^1\times S^1$
mają ściągalne przestrzenie
całkowite. Kontrakcja $v\mapsto(1-t)v$ ustala $0$,
więc każde oparte odwzorowanie sfery do $\R$ lub $\R^2$
staje się przez złożenie z nią homotopijne do stałego.
Stwierdzenie~\ref{prop:pi-n-nakrycie} daje zatem dla $n\geq2$
\[
 \pi_n(S^1)=0,\qquad
 \pi_n(T^2)=0.
\]
Jednocześnie
$H_2(T^2)=\Z$ z
Przykładu~\ref{prz:homologia-cw-torus}.
Już tu widać, że w
przestrzeni z nietrywialną
$\pi_1$ wyższa homologia
nie musi być obrazem
wyższej grupy homotopii.
\end{przyklad}


\section{Pierwsze odwzorowanie Hurewicza: od pętli do cykli}
\label{sec:hurewicz-jeden}

\begin{twierdzenie}[Hurewicz w stopniu pierwszym]
\label{thm:hurewicz-h1}
Dla przestrzeni łukowo spójnej $X$ odwzorowanie przypisujące
opartej pętli jej klasę singularnego $1$-cyklu indukuje naturalny
izomorfizm
\[
 h_1:\pi_1(X,x_0)_{\mathrm{ab}}
 =\pi_1(X,x_0)/[\pi_1(X,x_0),\pi_1(X,x_0)]
 \xrightarrow{\ \cong\ }H_1(X;\Z).
\]
\end{twierdzenie}
\begin{proof}
\emph{Krok 1: od konkatenacji dróg do sumy łańcuchów.}
Pętla $\ell:I\to X$, potraktowana jako singularny $1$-sympleks,
jest cyklem, ponieważ $\partial\ell=[x_0]-[x_0]=0$.
Potrzebujemy jednak sprawdzić, co dzieje się z działaniem grupowym:
łańcuch $\ell_1*\ell_2$ jest jednym generatorem $C_1(X)$,
a $\ell_1+\ell_2$ jest sumą dwóch generatorów. Nie są to te same łańcuchy.

Niech ogólniej $\alpha$ będzie drogą od $x$ do $y$, a $\beta$
drogą od $y$ do $z$. Dla współrzędnych barycentrycznych
$(t_0,t_1,t_2)$ trójkąta $\Delta^2$ połóżmy
$q(t_0,t_1,t_2)=t_1/2+t_2$ oraz
$T_{\alpha,\beta}=(\alpha*\beta)\circ q$.
Ten wybór wynika z czasów przejścia przez wierzchołki:
$q(v_0)=0$, $q(v_1)=1/2$, $q(v_2)=1$.
Afiniczne przedłużenie ma więc wartość
$t_0\cdot0+t_1\cdot(1/2)+t_2\cdot1$.
Na krawędzi $01$ parametr $q$ biegnie od $0$ do $1/2$,
więc ograniczeniem $T_{\alpha,\beta}$ jest dokładnie $\alpha$.
Na krawędzi $12$ otrzymujemy $\beta$, a na krawędzi $02$
całą drogę $\alpha*\beta$. Ze znaków w definicji brzegu wynika zatem
\[
 \partial T_{\alpha,\beta}
 =\beta-(\alpha*\beta)+\alpha.
\]
W szczególności konkatenacja pętli i suma ich łańcuchów
różnią się o brzeg. Homotopia oparta pętli także daje homologiczne
cykle przez operator pryzmatyczny z
Twierdzenia~\ref{thm:homotopia-singularna}.
Otrzymujemy więc dobrze określony homomorfizm
$\widehat h_1:\pi_1(X,x_0)\to H_1(X;\Z)$.
Ponieważ grupa docelowa jest abelowa, każdy komutator ma obraz zerowy:
obrazy $aba^{-1}b^{-1}$ sumują się do
$\widehat h_1(a)+\widehat h_1(b)-\widehat h_1(a)-\widehat h_1(b)=0$.
Stąd $\widehat h_1$ przechodzi na homomorfizm $h_1$
z abelianizacji podany w tezie.

\emph{Krok 2: odwzorowanie w przeciwnym kierunku.}
Dla każdego $x\in X$ wybierzmy drogę $u_x$ od $x_0$ do $x$;
istnieje ona dzięki spójności łukowej. Przyjmijmy, że $u_{x_0}$
jest drogą stałą. Singularnej krawędzi $\sigma:x\to y$ przypiszmy
klasę opartej pętli
$L(\sigma)=[u_x*\sigma*\bar u_y]$
w $\pi_1(X,x_0)_{\mathrm{ab}}$ i przedłużmy $L$ liniowo
na $C_1(X)$. Kreska nad drogą oznacza zmianę jej kierunku.

Sprawdźmy, że $L$ zeruje brzegi $2$-sympleksów.
Niech $\tau:\Delta^2\to X$ ma wierzchołki $a,b,c$ i krawędzie
$\tau_{01},\tau_{12},\tau_{02}$, skierowane zgodnie z indeksami.
W iloczynie pętli odpowiadających $\tau_{01}$ i $\tau_{12}$
środkowe drogi $\bar u_b*u_b$ można ściągnąć względem końców.
Po dopisaniu odwrotności pętli odpowiadającej $\tau_{02}$
tak samo skracają się drogi przy $c$. Pozostaje pętla
\[
 u_a*(\tau_{01}*\tau_{12}*\bar\tau_{02})*\bar u_a.
\]
Środkowa pętla jest brzegiem trójkąta odwzorowanego przez $\tau$,
więc jest ściągalna; doprowadzenie jej do bazy drogą $u_a$
nie zmienia tego faktu. W abelianizacji otrzymujemy zatem
\[
 L(\partial\tau)
 =L(\tau_{12})-L(\tau_{02})+L(\tau_{01})=0.
\]
Z liniowości $L$ zeruje całą grupę $\operatorname{im}\partial_2$.
Jego ograniczenie do cykli wyznacza więc homomorfizm
$\bar L:H_1(X;\Z)\to\pi_1(X,x_0)_{\mathrm{ab}}$.

\emph{Krok 3: sprawdzenie obu złożeń.}
Dla pętli opartej w $x_0$ obie pomocnicze drogi są stałe,
więc $u_{x_0}*\ell*\bar u_{x_0}$ jest homotopijna do $\ell$.
Stąd $(\bar L\circ h_1)([\ell])=[\ell]$.

Drugie złożenie wymaga rachunku na łańcuchach, ponieważ pojedyncze
$u_x$ zwykle nie są cyklami. Piszmy $\equiv$ dla równości
w $C_1(X)/\operatorname{im}\partial_2$.
Krok 1 daje $\alpha*\beta\equiv\alpha+\beta$ również dla dróg
o różnych końcach. Ponadto $\bar\alpha\equiv-\alpha$:
droga $\alpha*\bar\alpha$ jest opartą pętlą ściągalną,
zatem jest homologiczna do stałego $1$-sympleksu, a ten jest
brzegiem stałego $2$-sympleksu. Wobec tego
\[
 u_x*\sigma*\bar u_y\equiv u_x+\sigma-u_y.
\]
Niech $z=\sum_\sigma m_\sigma\sigma$ będzie cyklem.
Po zsumowaniu poprzednich równości łańcuch reprezentujący
$h_1(\bar L([z]))$ jest równy modulo brzegi
\[
 z+\sum_{v\in X}
 \left(\sum_{\sigma(0)=v}m_\sigma
       -\sum_{\sigma(1)=v}m_\sigma\right)u_v.
\]
Wszystkie sumy są skończone. Wyrażenie w nawiasie jest
przeciwieństwem współczynnika przy $[v]$ w
$\partial z=\sum_\sigma m_\sigma([\sigma(1)]-[\sigma(0)])$.
Ponieważ $\partial z=0$ i punkty $[v]$ są niezależnymi generatorami
$C_0(X)$, każdy z tych współczynników jest zerowy.
Pozostaje dokładnie $z$, co dowodzi $h_1\circ\bar L=\id$.
Oba homomorfizmy są więc odwrotne. W szczególności wynik nie zależy
od pomocniczych dróg $u_x$: wszystkie ich wybory dają odwrotność
tego samego $h_1$.

\emph{Krok 4: naturalność.}
Dla odwzorowania opartego $f:(X,x_0)\to(Y,y_0)$ obraz pętli
na obu drogach porównania jest tą samą pętlą $f\circ\ell$.
Zatem $h_1(f_*[\ell])=f_*(h_1[\ell])$.
Równość przechodzi na abelianizacje, bo homomorfizm grup
przenosi komutatory w komutatory. To dowodzi naturalności.
\end{proof}

\begin{przyklad}[Od wolnej grupy do liczby pętli]
Bukiet $r$ okręgów ma grupę podstawową $F_r$
z Przykładu~\ref{prz:pi1-dwa-okregi}; dla $r=2$ rachunek
podano tam wprost. Oznaczmy standardowe pętle przez
$a_1,\ldots,a_r$. W abelianizacji możemy przestawiać litery,
więc każde słowo ma postać $a_1^{m_1}\cdots a_r^{m_r}$.
Wykładnik $m_i$ jest sumą wykładników litery $a_i$ w słowie:
skracanie sąsiednich odwrotnych liter go nie zmienia.
Odwzorowanie na $(m_1,\ldots,m_r)\in\Z^r$ jest więc dobrze
określone, a jego odwrotność przypisuje wektorowi klasę
$a_1^{m_1}\cdots a_r^{m_r}$ w abelianizacji.
Stąd $F_r^{\mathrm{ab}}\cong\Z^r$, zgodnie z
Przykładem~\ref{prz:homologia-bukiet}.
Dla dwóch okręgów pętla $a_1*a_2*\bar a_1*\bar a_2$
jest nietrywialnym słowem zredukowanym w $F_2$, ale w $H_1$
ma klasę $(1,0)+(0,1)-(1,0)-(0,1)=(0,0)$.
Homologia pamięta tu liczby obiegów, lecz gubi kolejność
przebywania pętli.
\end{przyklad}

\section{Twierdzenie Whiteheada}
\label{sec:whitehead}

W tym twierdzeniu istotne jest konkretne odwzorowanie
$f:X\to Y$ i homomorfizmy, które ono indukuje.
Samo istnienie abstrakcyjnych izomorfizmów między wszystkimi
grupami $\pi_n(X)$ i $\pi_n(Y)$ nie jest jego założeniem.
Dla kompleksów CW izomorfizmy indukowane przez $f$ pozwolą
przechodzić od deformacji sfer do deformacji całych komórek,
a następnie skonstruować homotopijną odwrotność $f$.

\begin{definicja}[Słaba równoważność homotopijna]
\index{slaba rownowaznosc homotopijna@Słaba równoważność homotopijna}
\label{def:slaba-rownowaznosc}
Odwzorowanie $f:X\to Y$ jest \emph{słabą równoważnością homotopijną},
jeśli daje bijekcję $\pi_0(X)\to\pi_0(Y)$ i dla wszystkich punktów
bazowych $x\in X$ oraz wszystkich $n\geq1$ izomorfizmy
$f_*:\pi_n(X,x)\to\pi_n(Y,f(x))$.
\end{definicja}

Każda równoważność homotopijna jest słabą równoważnością:
homotopijna odwrotność daje odwrotne homomorfizmy po uwzględnieniu
zmiany punktu bazowego opisanej po
Stwierdzeniu~\ref{prop:pi-n-funktor}.
Whitehead dowodzi implikacji odwrotnej dla CW.
Potrzebny krok geometryczny zawiera następujący lemat.

\begin{lemat}[Ściskanie komórek]
\label{lem:kompresja-cw}
Niech $(K,L)$ będzie parą CW, a $B\subseteq Y$ niepustą
podprzestrzenią. Załóżmy, że każda składowa łukowa $Y$ spotyka $B$
oraz wszystkie grupy $\pi_k(Y,B,b)$ dla $k\geq1$ i $b\in B$
są trywialne; dla $k=1$ chodzi o zbiór z wyróżnioną klasą.
Wówczas każde odwzorowanie par $F:(K,L)\to(Y,B)$ jest homotopijne
\emph{względem $L$} z odwzorowaniem o obrazie w $B$.
\end{lemat}
\begin{proof}
\emph{Krok 1: wierzchołki i założenie indukcyjne.}
Każdy obraz $0$-komórki spoza $L$ łączymy drogą z punktem $B$;
jest to możliwe dzięki założeniu o składowych $Y$.
Na wierzchołkach z $L$ wybieramy homotopie stałe.
Lemat~\ref{lem:hep-cw} przedłuża je z podkompleksu $L\cup K^0$
na całe $K$. Końcowe odwzorowanie wysyła $L\cup K^0$ do $B$.

Załóżmy indukcyjnie, że bieżące odwzorowanie wysyła
$L\cup K^{k-1}$ do $B$. Na tym podkompleksie nie będziemy
już zmieniać wartości. Jeśli $\chi:D^k\to K$ jest odwzorowaniem
charakterystycznym komórki spoza $L$, to $f=F\circ\chi$
wysyła cały $S^{k-1}$ do $B$. Wybieramy jeden punkt tego brzegu
jako punkt bazowy; jego obraz to pewne $b\in B$.
W ten sposób $f$ reprezentuje klasę w $\pi_k(Y,B,b)$.
Jej zerowość daje homotopię $G:D^k\times I\to Y$ od $f$
do stałego odwzorowania o wartości $b$, przy czym
$G_t(S^{k-1})\subseteq B$. Dla $k=1$ jest to właśnie
zerowość względnej klasy drogi, z jednym końcem unieruchomionym w $b$.

\emph{Krok 2: zatrzymanie całego brzegu.}
Homotopia $G$ może przesuwać pozostałe punkty brzegu w $B$.
Nie można jeszcze skleić takich homotopii sąsiednich komórek,
bo na wspólnym brzegu mogłyby dawać różne wartości.
Zachowamy więc na zewnętrznym brzegu początkowe $f$,
a przesuwający się brzeg schowamy w kołnierzu dysku.
Wybierzmy wewnętrzny promień $a_t=1-t/2$:
kołnierz ma grubość $t/2$ i zanika w chwili $0$,
a wewnętrzny dysk nigdy nie zapada się do punktu.
W jego wnętrzu dzielimy promień przez $a_t$, aby ponownie
uzyskać dysk jednostkowy. Na kołnierzu czas homotopii
ma maleć od $t$ przy $r=a_t$ do $0$ przy $r=1$.
Liniowa interpolacja dla $t>0$ daje
$t(1-r)/(1-a_t)=2(1-r)$; ostatni wzór ma sens także przy $t=0$.
Dla $x=ru$, gdzie $0\leq r\leq1$
i $u\in S^{k-1}$, definiujemy
\[
 \widehat G_t(ru)=
 \begin{cases}
 G_t(ru/a_t),&0\leq r\leq a_t,\\
 G_{2(1-r)}(u),&a_t\leq r\leq1.
 \end{cases}
\]
W środku dysku używamy tylko pierwszego wzoru,
który jest niezależny od nieokreślonego tam kierunku $u$.
Na styku $r=a_t$ mamy $2(1-r)=t$, więc oba wzory dają $G_t(u)$.
Na zewnętrznym brzegu $r=1$ wartość wynosi stale $G_0(u)=f(u)$.
Dla $t=0$ pierwszy wzór daje $f$ na całym dysku.
Choć kołnierz zanika przy $t=0$, ciągłość pozostaje zachowana:
na nim czas $2(1-r)$ leży w $[0,t]$, więc przy $t\to0$
i zbliżaniu się do brzegu wartości obu wzorów dążą do tej samej
wartości $G_0$.

Dla $t=1$ wewnętrzny dysk ma stały obraz $b$,
a kołnierz korzysta wyłącznie z wartości $G_s(u)$ dla $u\in S^{k-1}$,
które leżą w $B$. Otrzymaliśmy więc deformację $f$ do $B$
\emph{nieruchomą na całym brzegu}.
Wykonujemy tę konstrukcję na wszystkich $k$-komórkach spoza $L$.
Na ich brzegach homotopie są stałe, zatem sklejają się
z homotopią stałą na $L\cup K^{k-1}$.
Lemat~\ref{lem:hep-cw} przedłuża je z $L\cup K^k$ na całe $K$.
Końcowe odwzorowanie wysyła już $L\cup K^k$ do $B$.

\emph{Krok 3: sklejenie nieskończenie wielu etapów.}
Jeśli wymiar $K$ jest skończony, po skończenie wielu krokach
otrzymujemy tezę. Dla dowolnego CW wykonujemy etap dotyczący
$k$-komórek na przedziale
$[1-2^{-k},1-2^{-(k+1)}]$, $k=0,1,\ldots$,
i w chwili $1$ bierzemy ustaloną końcową wartość w każdym punkcie.
Po etapie $k$ cały $k$-szkielet pozostaje nieruchomy,
ponieważ późniejsze homotopie są względem niego stałe.
Na domkniętej komórce wymiaru $d$ występuje więc tylko
skończenie wiele niestałych etapów; od etapu $d$ do chwili $1$
odwzorowanie już się tam nie zmienia.
Złożenie z $\chi\times\id_I:D^d\times I\to K\times I$
jest zatem ciągłe także przy czasie $1$.
Kryterium z Lematu~\ref{lem:zwartosc-cw-pi} daje ciągłość
homotopii na całym $K\times I$. Na $L$ jest ona przez cały czas
stała, a końcowy obraz każdej komórki leży w $B$.
\end{proof}

\begin{twierdzenie}[Whitehead]
\label{thm:whitehead}
Słaba równoważność homotopijna $f:X\to Y$ między kompleksami CW
jest równoważnością homotopijną. Jeśli $i:X\hookrightarrow Y$
jest inkluzją podkompleksu i indukuje izomorfizmy wszystkich
$\pi_n$ oraz bijekcję na $\pi_0$, to $X$ jest silnym retraktem
deformacyjnym $Y$.
\end{twierdzenie}
\begin{proof}
Jeśli $X=\varnothing$, bijekcja na składowych wymusza
$Y=\varnothing$ i teza jest natychmiastowa.
Dalej zakładamy $X\ne\varnothing$.

\emph{Krok 1: zanik grup względnych dla inkluzji.}
Niech najpierw $i:X\hookrightarrow Y$ będzie inkluzją
podkompleksu z tezy. Dla każdej bazy $x_0\in X$
w długim ciągu~\eqref{eq:les-pi} występuje fragment
\[
 \pi_n(X,x_0)\xrightarrow[\cong]{i_*}\pi_n(Y,x_0)
 \xrightarrow{j_*}\pi_n(Y,X,x_0)
 \xrightarrow{\partial}\pi_{n-1}(X,x_0)
 \xrightarrow[\cong]{i_*}\pi_{n-1}(Y,x_0).
\]
Dla $n\geq2$ weźmy dowolną klasę $a\in\pi_n(Y,X,x_0)$.
Jej brzeg leży w jądrze ostatniego $i_*$, więc $\partial a=0$.
Dokładność daje $a=j_*(b)$ dla pewnego $b\in\pi_n(Y,x_0)$.
Surjektywność pierwszego $i_*$ pozwala napisać $b=i_*(c)$,
a wtedy $a=(j_*\circ i_*)(c)=0$.
Ten argument działa również dla $n=2$: nie wymaga przemienności
$\pi_2(Y,X,x_0)$.

Stopień $1$ sprawdzamy osobno, bo występuje w nim względny
zbiór z wyróżnioną klasą. Reprezentuje go droga $\gamma$
od pewnego $a\in X$ do ustalonego $x_0\in X$.
Jej końce leżą w tej samej składowej łukowej $Y$.
Injektywność $\pi_0(X)\to\pi_0(Y)$ oznacza, że leżą również
w tej samej składowej $X$; wybieramy więc drogę
$\lambda:x_0\to a$ w $X$.
Pętla $\lambda*\gamma$ w $Y$ jest, dzięki surjektywności
$\pi_1(X,x_0)\to\pi_1(Y,x_0)$, homotopijna względem końców
z pewną pętlą $\mu$ w $X$.
Po dopięciu $\bar\lambda$ i skróceniu $\bar\lambda*\lambda$
otrzymujemy homotopię $\gamma\simeq\bar\lambda*\mu$
względem obu końców. Prawa strona jest drogą w $X$.
Każdą drogę $\rho:I\to X$ z $\rho(1)=x_0$ można ściągnąć
względnie $X$ wzorem $\rho_t(s)=\rho((1-t)s+t)$:
koniec $s=1$ pozostaje w $x_0$, a drugi koniec porusza się w $X$.
Zatem $\pi_1(Y,X,x_0)$ ma tylko wyróżnioną klasę.
Każda składowa $Y$ spotyka $X$, bo odwzorowanie na $\pi_0$
jest także surjekcją.

\emph{Krok 2: silna retrakcja dla inkluzji.}
Możemy zastosować Lemat~\ref{lem:kompresja-cw} do
$F=\id_Y:(Y,X)\to(Y,X)$: wszystkie jego założenia
sprawdziliśmy w Kroku 1 dla każdej bazy w $X$.
Otrzymujemy homotopię od $\id_Y$ do odwzorowania $s:Y\to X$,
która na $X$ jest stale identycznością.
W szczególności $s|_X=\id_X$, a homotopia jest
silną retrakcją deformacyjną $Y$ na $X$.

\emph{Krok 3: zamiana dowolnego odwzorowania na inkluzję.}
Niech teraz $f:X\to Y$ będzie dowolną słabą równoważnością CW.
Twierdzenie~\ref{thm:przyblizenie-komorkowe} daje komórkowe
$f'\simeq f$. Homotopia może przesuwać obraz punktu bazowego,
lecz zmiana bazy wzdłuż tej drogi utożsamia indukowane
homomorfizmy; stąd $f'$ nadal jest słabą równoważnością.
Utwórzmy cylinder odwzorowania
\[
 M_{f'}=(X\times I)\sqcup Y/
       \bigl((x,1)\sim f'(x)\bigr).
\]
Poza komórkami $Y$ ma on dolne komórki $e^k\times\{0\}$
i komórki $e^k\times(0,1)$ wymiaru $k+1$.
Brzeg charakterystycznego walca $D^k\times I$ składa się
z dołu, boków nad $\partial D^k$ i góry odwzorowanej przez
$f'$ do $Y^k$. Wszystkie te części leżą w $k$-szkielecie
konstruowanego cylindra: tutaj korzystamy z komórkowości $f'$.
Warunek skończoności domknięć wynika z niego dla $X,Y$
oraz ze zwartości obrazów dysków charakterystycznych.
Topologia ilorazowa wraz z kryterium z
Lematu~\ref{lem:zwartosc-cw-pi} daje topologię słabą tej struktury.
Zatem $M_{f'}$ jest CW, a dolna inkluzja
$j:X\hookrightarrow M_{f'}$, $j(x)=[x,0]$, jest inkluzją podkompleksu.

Oznaczmy inkluzję drugiej części przez $i:Y\hookrightarrow M_{f'}$.
Cylinder ma retrakcję $r:M_{f'}\to Y$ daną przez
$r([x,u])=f'(x)$ oraz $r([y])=y$.
Homotopia $D_t([x,u])=[x,u+t(1-u)]$, $D_t([y])=[y]$,
przesuwa każdy odcinek do jego górnego końca i pozostawia $Y$
nieruchome. Wzory są zgodne na utożsamionych punktach
$[x,1]=[f'(x)]$, więc wyznaczają ciągłą homotopię na ilorazie;
ciągłość z parametrem wynika z tego samego ilorazowego kryterium.
Otrzymujemy $r\circ i=\id_Y$ oraz $i\circ r\simeq\id_{M_{f'}}$.
Ponadto $r\circ j=f'$.

Ponieważ $r$ jest równoważnością homotopijną, indukuje
izomorfizmy grup homotopii i bijekcję na składowych.
Równość $r_*\circ j_*=f'_*$ pokazuje, że także $j_*$
jest izomorfizmem w każdym stopniu: jest złożeniem
$(r_*)^{-1}\circ f'_*$, z bazami $j(x)$ i $f'(x)$.
Tak samo uzyskujemy bijekcję na $\pi_0$.
Kroki 1--2 stosujemy więc do inkluzji $j$.
Dają retrakcję $s:M_{f'}\to X$ i relacje
$s\circ j=\id_X$, $j\circ s\simeq\id_{M_{f'}}$.

\emph{Krok 4: sprawdzenie odwrotności.}
Połóżmy $g=s\circ i:Y\to X$.
Korzystając najpierw z homotopii $i\circ r\simeq\id_{M_{f'}}$,
a potem z $j\circ s\simeq\id_{M_{f'}}$, obliczamy
\[
 \begin{aligned}
 g\circ f'&=s\circ i\circ r\circ j
       \simeq s\circ j=\id_X,\\
 f'\circ g&=r\circ j\circ s\circ i
       \simeq r\circ i=\id_Y.
 \end{aligned}
\]
Homotopie można składać z $s,j$ lub $r,i$, bo złożenie
ciągłych odwzorowań z homotopią znów jest homotopią.
Wreszcie $f\simeq f'$ daje $g\circ f\simeq g\circ f'$
oraz $f\circ g\simeq f'\circ g$.
To samo $g$ jest więc homotopijną odwrotnością $f$.
\end{proof}

\begin{figure}[H]
\centering
\begin{tikzpicture}[>=Stealth,line cap=round,line join=round]
 \fill[manifoldblue!10]
  (0,0)--(4.1,0)--(4.35,2.0)--(.25,2.0)--cycle;
 \draw[manifoldblue!70!black,thick]
  (0,0)--(4.1,0)--(4.35,2.0)--(.25,2.0)--cycle;
 \foreach \x in {.7,1.8,2.9,3.7}{
  \draw[manifoldblue!35] (\x,.07)--({\x+.23},1.94);
 }
 \fill[manifoldgreen!18]
  (-.38,2.0) rectangle (4.86,2.5);
 \draw[manifoldgreen!70!black,thick]
  (-.38,2.0) rectangle (4.86,2.5);
 \node at (2.2,2.25) {$Y$};
 \node[below] at (2.08,0) {$X\times\{0\}$};
 \node[manifoldblue!70!black] at (2.2,1.02)
  {$X\times I$};
 \node[manifoldred!85!black,above]
  at (2.2,2.62) {$(x,1)\sim f'(x)$};
 \draw[->,manifoldred,thick]
  (3.35,1.48)--(3.35,2.2)
  node[midway,right=3pt] {$r$};
 \node[font=\small,align=left,anchor=west,text width=4.8cm]
  at (5.15,1.15)
  {Cylinder zwija się na $Y$.
   Gdy $f'$ jest słabą równoważnością,
   poprzednia część dowodu zwija go
   również na dolną kopię $X$.};
\end{tikzpicture}
\caption{Cylinder odwzorowania zamienia
odwzorowanie $f':X\to Y$ na inkluzję
podkompleksu $X\hookrightarrow M_{f'}$.
Rysunek pokazuje mechanizm dowodu
twierdzenia Whiteheada.}
\label{fig:cylinder-whiteheada}
\end{figure}

\begin{uwaga}[Zakres twierdzenia]
Sprawdzenie tylko $\pi_1$ i kilku pierwszych grup $\pi_n$
nie wystarcza do zastosowania Whiteheada: założenie dotyczy
\emph{każdego} $n$ i homomorfizmów indukowanych przez to samo $f$.
Również samo podobieństwo grup homologii dwóch przestrzeni
nie wskazuje jeszcze odwzorowania, które mogłoby być równoważnością.
Niżej wyprowadzimy użyteczne kryterium homologiczne
przy dodatkowej jednospójności.
\end{uwaga}

\section{Odwzorowanie Hurewicza i pierwszy niezerowy wymiar}
\label{sec:hurewicz-wyzszy}

\begin{definicja}[Odwzorowanie Hurewicza]
\index{odwzorowanie hurewicza@Odwzorowanie Hurewicza}
\label{def:mapa-hurewicza}
Homologia $H_n(S^n;\Z)$
jest izomorficzna z $\Z$
na mocy
Twierdzenia~\ref{thm:homologia-sfer}.
Wybierzmy generatory zgodnie z orientacjami dysków i ich brzegów,
tak aby $\delta[D^n,S^{n-1}]=[S^{n-1}]$.
Dla $n\geq2$ określamy
\[
 h_n:\pi_n(X,x_0)\longrightarrow H_n(X;\Z),
 \qquad h_n([f])=f_*[S^n].
\]
Dla pary $(X,A)$ bierzemy
analogicznie generator
$[D^n,S^{n-1}]
\in H_n(D^n,S^{n-1})$ i
kładziemy
\[
 h_n^{\mathrm{rel}}([F])
 =F_*[D^n,S^{n-1}]
 \in H_n(X,A).
\]
W stopniu $1$ dla przestrzeni używamy homomorfizmu
$\pi_1(X)\to\pi_1(X)_{\mathrm{ab}}\xrightarrow{h_1}H_1(X)$.
Względny homomorfizm określamy tutaj dla $n\geq2$;
zbiór $\pi_1(X,A)$ nie jest w ogólności grupą.
\end{definicja}

Homotopie oparte są w
szczególności zwykłymi
homotopiami odwzorowań, więc
$h_n$ jest dobrze
określone dzięki
Twierdzeniu~\ref{thm:homotopia-singularna}.
Odwzorowanie jest homomorfizmem:
sklejanie dwóch połówek
sfery odpowiada sumie
ich klas fundamentalnych.
Dokładniej, sklejanie jest złożeniem odwzorowania ściskającego równik
$S^n\to S^n\vee S^n$ z $f\vee g$. Odwzorowanie ściskające posyła
generator na sumę dwóch generatorów: rzut na każdy składnik
ma stopień $+1$, co widać z klasy orientacji odpowiedniej półkuli
i naturalności odwzorowania łączącego dla dysku i brzegu.
Względnie używamy kostki podzielonej na dwie połowy w pierwszej współrzędnej.
Po zorientowaniu obu połówek zgodnie z całą kostką ich wewnętrzne
ściany kasują się; obraz łańcucha fundamentalnego jest sumą obu obrazów.
To daje również homomorfizm względny. Naturalność wynika od razu ze składania odwzorowań.

Odwzorowania Hurewicza są także zgodne ze strzałkami brzegowymi
obu długich ciągów pary. Oznaczmy je chwilowo przez
$\partial_\pi$ w homotopii i $\delta_H$ w homologii.
Dla reprezentanta $F:(D^n,S^{n-1},*)\to(X,A,x_0)$ naturalność
ciągu homologii daje
\[
 \begin{aligned}
 \delta_H\bigl(h_n^{\mathrm{rel}}[F]\bigr)
 &=\delta_H\bigl(F_*[D^n,S^{n-1}]\bigr)\\
 &=(F|_{S^{n-1}})_*\bigl(\delta_H[D^n,S^{n-1}]\bigr)\\
 &=(F|_{S^{n-1}})_*[S^{n-1}]
 =h_{n-1}\bigl(\partial_\pi[F]\bigr).
 \end{aligned}
\]
Dla $n=2$ ostatnia strzałka jest homomorfizmem
$\pi_1(A)\to H_1(A)$ z Twierdzenia~\ref{thm:hurewicz-h1}.
Uzgodniony wybór orientacji dysku i jego brzegu zapewnia znak dodatni.

Do dowodu potrzebujemy
przeliczyć grupy homotopii
prostych modeli komórkowych.
Najważniejszy jest
następujący krok: w
pierwszym nowym wymiarze
komórki są niezależne,
a komórki o wymiar
wyższym wprowadzają
relacje.

\begin{lemat}[Pierwsze komórki i ich relacje]
\label{lem:pierwsze-komorki-pi}
Niech $A$ będzie jednospójnym
podkompleksem CW kompleksu $Z$.
Załóżmy, że wszystkie komórki
$Z\setminus A$ mają wymiar
co najmniej $n\geq2$.
Oznaczmy przez $I_n$
zbiór nowych $n$-komórek.
Wtedy
\[
 \pi_n(Z^n\cup A,A)
 \cong\bigoplus_{\alpha\in I_n}\Z[e_\alpha^n]
 \xrightarrow[\cong]{\ h_n^{\mathrm{rel}}\ }
 H_n(Z^n\cup A,A).
\]
Po dołączeniu $(n+1)$-komórek
obie grupy otrzymujemy
z tej wolnej grupy przez
te \emph{same} relacje:
obraz brzegu odwzorowania
dołączającego $S^n\to
Z^n\cup A$. Komórki
o wymiarach co najmniej
$n+2$ nie zmieniają
tych grup.
\end{lemat}
\begin{proof}
Niech $W=Z^n\cup A$ oraz $W'=Z^{n+1}\cup A$.
Przestrzeń $W$ jest jednospójna: do jednospójnego $A$
dołączamy tylko komórki wymiaru co najmniej $2$, co nie tworzy
nowych pętli (lokalny krok przybliżenia komórkowego przesuwa pętlę do $A$).
W szczególności także $\pi_2(W,A)$ jest abelowa, bo z ciągu pary
jest ilorazem abelowej $\pi_2(W)$ przy $\pi_1(A)=0$.
Możemy więc mówić o homomorfizmie z wolnej grupy abelowej
również w przypadku $n=2$.
Po zwinięciu $A$
otrzymujemy bukiet
$\bigvee_{\alpha\in I_n}S^n$.
Względna homologia
$H_n(W,A)$ jest
wolna na komórkach
na mocy
Twierdzenia~\ref{thm:warstwa-komorek}
i rozszerzenia na pary CW po Lemacie~\ref{lem:zwartosc-cw-pi};
to jeszcze nie jest
argument o homotopii,
więc rozpiszmy go
osobno.

Każde odwzorowanie charakterystyczne nowej komórki opieramy w $x_0$
przez dołączenie drogi w $A$ do obrazu wybranego punktu jej brzegu.
Wybór drogi nie zmienia klasy, bo $A$ jest jednospójna.
Tak rozumiemy symbole $[e_\alpha^n]$ po stronie homotopii.
Weźmy teraz odwzorowanie względne
$F:(D^n,S^{n-1},*)\to(W,A,x_0)$.
Zwarty obraz spotyka
tylko skończenie wiele
otwartych komórek.
W każdej z nowych $n$-komórek stosujemy przybliżenie
z Uwagi~\ref{rem:przyblizenie-rowne-wymiary}, nieruchome przy $S^{n-1}$.
Wybieramy regularny punkt $z_\alpha$ poza obrazami ścian triangulacji.
Ma on skończenie wiele przeciwobrazów; na małych otoczeniach
$F$ jest afinicznym homeomorfizmem o znaku $+1$ lub $-1$.
Wybierzmy wokół $z_\alpha$ małą kulę. Jej przeciwobrazy są rozłącznymi
dyskami $D_j$, a poza nimi $F$ omija tę kulę.
Parametryzując $D_j$ przez odwrotność lokalnego odwzorowania afinicznego,
otrzymujemy dokładnie lokalny fragment odwzorowania charakterystycznego;
ujemny znak odpowiada odbiciu pierwszej współrzędnej, czyli odwrotności w grupie.

Rozpiszmy jednoczesne rozszerzenie małych dysków i przesunięcie
ich dopełnienia do $A$. We współrzędnych jednej komórki niech
$\overline B_\varepsilon(z)\subset\operatorname{int}D^n$ będzie
wybraną małą kulą. Dla jednostkowego $u$ oznaczmy przez $R(u)>\varepsilon$
odległość od $z$ do jedynego punktu przecięcia promienia $z+ru$
z $\partial D^n$. Jest to funkcja $R_z(u)$ obliczona
w dowodzie Twierdzenia~\ref{thm:przyblizenie-komorkowe}.
Aby rozciągnąć promień małej kuli $[0,\varepsilon]$
na $[0,R(u)]$, mnożymy $r$ przez $R(u)/\varepsilon$.
Interpolacja od identyczności daje
$(1-t)r+trR(u)/\varepsilon$. Dla większych $r$
wartość mogłaby wyjść poza dysk, więc zatrzymujemy ją
na $R(u)$, biorąc minimum. Na całym dysku stosujemy homotopię
\[
 \Psi_t(z+ru)=z+
 \min\!\left\{(1-t)r+t\frac{rR(u)}{\varepsilon},\ R(u)\right\}u,
 \qquad \Psi_t(z)=z.
\]
Funkcja $R$ jest ciągła i ograniczona przez $1+\|z\|$.
Odległość obrazu od $z$ jest nie większa niż
$r(1+(1+\|z\|)/\varepsilon)$, niezależnie od $t$ i $u$,
więc wzór jest ciągły także w $z$. Dla $t=0$ mamy identyczność, a punkty brzegu
$r=R(u)$ pozostają nieruchome. Dla $t=1$ cała mała kula
rozciąga się homeomorficznie na $D^n$, podczas gdy jej dopełnienie
trafia do $\partial D^n$. Ponieważ na brzegu niczego nie zmieniamy,
te homotopie schodzą przez odwzorowania charakterystyczne na $W$
i są identycznością na $A$. Po złożeniu z $F$ otrzymujemy
na każdym $D_j$ cały dysk charakterystyczny, a poza wnętrzami
tych dysków odwzorowanie o obrazie w $A$.

Pozostaje uzasadnić, dlaczego takie odwzorowanie jest sumą klas
lokalnych dysków. W dopełnieniu ich wnętrz wybieramy skończone
drzewo łuków łączących je z punktem bazowym. Łuki można dobrać
odcinkowo liniowo, z rozłącznymi wnętrzami; w wymiarze dwa
prowadzimy je kolejno w dysku z otworami. Pogrubiamy łuki do
cienkich pasów. Pojedynczy pas doprowadza jeden dysk do zewnętrznego
brzegu. Rozcięcie wzdłuż niego oddziela tę część od pozostałych;
po przeparametryzowaniu obu części są one dwiema połówkami kostki
z Definicji~\ref{def:pi-wzgledne}. Aby użyć tego dodawania,
musimy jeszcze uzyskać stałą wartość $x_0$ na ścianie rozcięcia.
Niech $E\cong D^{n-1}$ będzie tą ścianą, zawierającą wyróżniony
punkt $p\in\partial D^n$, z $F(p)=x_0$. Cięcia i pasy wybieramy
wielościennie, więc $E$, $\partial D^n$ i ich suma są podkompleksami
odpowiednio podzielonego dysku.
Obraz $F(E)$ leży w $A$. Kontrakcja $E$ do jego punktu $p$,
nieruchoma w $p$, daje po złożeniu z $F$ homotopię
$F|_E$ do stałego $x_0$, w całości w $A$.
Jej ograniczenie do $E\cap\partial D^n$ przedłużamy najpierw
na $\partial D^n$ z wartościami w $A$, korzystając z HEP.
Dostajemy zgodne dane homotopii na $E\cup\partial D^n$;
HEP dla pary $(D^n,E\cup\partial D^n)$ przedłuża je na cały dysk
z wartościami w $W$. Podczas tej deformacji brzeg pozostaje
w $A$, a $p$ w $x_0$, więc klasa względna się nie zmienia.
Teraz $E$ jest wysłane stale do $x_0$ i obie części rzeczywiście
dają kostkową sumę, po przeparametryzowaniu i zwężeniu stałych
fragmentów brzegu. Doprowadzenia wzdłuż pasów zapisują oba składniki
przy tej samej bazie $x_0$.
Powtarzamy tę operację, usuwając kolejne końcowe gałęzie drzewa.
Po skończenie wielu cięciach pozostają oddzielne dyski charakterystyczne
oraz części z całym obrazem w $A$. Te ostatnie dają zero względne:
ściągamy ich dziedzinę do wyróżnionego punktu, a wszystkie wartości
przez cały czas pozostają w $A$.

Dwie drogi doprowadzające ten sam dysk różnią się pętlą w $A$.
Jej nullhomotopia, wynikająca z $\pi_1(A)=0$, daje homotopię
obu pasów przy ustalonych końcach. Wybór drzewa i dróg nie zmienia
więc uzyskanych klas. Jeśli $\alpha(j)$ oznacza komórkę zawierającą
obraz lokalnego dysku $D_j$, a $\varepsilon_j\in\{1,-1\}$ jego znak,
otrzymujemy konkretną równość
$[F]=\sum_j\varepsilon_j[e_{\alpha(j)}^n]$.
Suma jest skończona dzięki zwartości dziedziny i skończonej liczbie
przeciwobrazów wybranych punktów.
To dowodzi surjektywności
odwzorowania od wolnej grupy
na nowych komórkach
do $\pi_n(W,A)$.
Obrazy tych generatorów
przez $h_n^{\mathrm{rel}}$
tworzą bazę
$H_n(W,A)$. Gdyby ich
kombinacja dawała
zerową klasę względnej
homotopii, jej obraz
homologiczny również
byłby zerowy; niezależność
bazy wymusza wszystkie
współczynniki równe $0$.
Mamy więc izomorfizm.

Niech teraz dołączamy dyski $D^{n+1}_\beta$, otrzymując $W'$.
Każde odwzorowanie $D^n\to W'$ względnie $A$ przesuwamy do $W$
przez lokalne usuwanie punktów w komórkach wymiaru $n+1$.
Inkluzja indukuje więc surjekcję na $\pi_n$.
Relację od $\varphi_\beta:S^n\to W$ rozumiemy jako jej obraz
przez $\pi_n(W)\to\pi_n(W,A)$.
Ich odwzorowania charakterystyczne
dają oczywiste relacje:
brzeg każdego takiego
dysku staje się zerem
w $\pi_n(W',A)$.
Dlaczego nie pojawiają się inne relacje?
Niech $F_0,F_1:(D^n,S^{n-1},*)\to(W,A,x_0)$ mają
ten sam obraz w $\pi_n(W',A)$ i niech
$H:D^n\times I\to W'$ będzie homotopią względną między nimi.
Zatem $H|_{D^n\times\{i\}}=F_i$,
$H(S^{n-1}\times I)\subseteq A$ oraz $H(*,t)=x_0$.
Możemy użyć kostkowego modelu $D^n=I^n$;
wtedy cylinder $B=D^n\times I$ jest wielościennym
$(n+1)$-dyskiem. Orientujemy go tak, aby górny dysk
miał orientację $D^n$, a dolny orientację przeciwną.
Cały brzeg $B$ ma obraz w $W$.

Obraz $H$ spotyka tylko skończenie wiele nowych komórek.
Uwaga~\ref{rem:przyblizenie-rowne-wymiary}, zastosowana
w wymiarze $n+1$ i względem całego brzegu $B$, pozwala
wybrać w każdej z tych komórek punkt $z_\beta$ tak, aby
wszystkie jego przeciwobrazy leżały we wnętrzach sympleksów,
na których $H$ jest afinicznym homeomorfizmem.
Przeciwobrazów jest skończenie wiele.
Wokół każdego wybieramy mały wypukły wielościenny dysk
$Q_j\subset\operatorname{int}B$, zawarty w takim sympleksie.
Dyski $Q_j$ wybieramy parami rozłączne; każdy zawiera
dokładnie jeden z rozpatrywanych przeciwobrazów.
Oznaczmy przez $\beta(j)$ odpowiadającą mu komórkę
i przez $\varepsilon_j\in\{1,-1\}$ znak lokalnego
odwzorowania afinicznego, względem wybranych orientacji.

Połóżmy
\[
 P=B\setminus\bigcup_j\operatorname{int}Q_j.
\]
Jest to skończony wielościan z brzegiem. Możemy triangulować
go bez korzystania z ogólnego twierdzenia o triangulacji
rozmaitości: dzielimy kostkę $B$ hiperpłaszczyznami
zawierającymi ściany wszystkich $Q_j$. Powstaje skończenie
wiele wypukłych wielościanów, które triangulujemy przez
barycentra ich ścian. W tym podziale $P$, $\partial B$
i wszystkie $\partial Q_j$ są podkompleksami.
Następnie wypełnienie każdego otworu traktujemy jako
dołączenie do $P$ jednej $(n+1)$-komórki $\operatorname{int}Q_j$.
Daje to parę CW $(B,P)$.

Przestrzeń $P$ jest łukowo spójna: fragment drogi w kostce,
który przechodzi przez $\operatorname{int}Q_j$, możemy
zastąpić drogą w łukowo spójnym brzegu $\partial Q_j$.
Jest również jednospójna. Wklejenie jednego $Q_j$ nie zmienia
grupy podstawowej: w twierdzeniu van Kampena używamy
otwartych otoczeń z kołnierzami, które retrahują się
odpowiednio na dotychczasową przestrzeń, $Q_j$ i ich
przecięcie $\partial Q_j\cong S^n$.
Ostatnia sfera jest jednospójna, bo $n\geq2$,
a sam dysk jest ściągalny. Inkluzja indukuje więc
izomorfizm grup podstawowych. Po wklejeniu wszystkich
dysków otrzymujemy ściągalne $B$, zatem $\pi_1(P)=0$.

Teraz stosujemy \emph{udowodnioną już pierwszą część}
niniejszego lematu w wymiarze $n+1$, z bazą $P$.
Otrzymujemy wolną grupę na klasach charakterystycznych
dysków $Q_j$ oraz izomorfizm
\[
 h_{n+1}^{\mathrm{rel}}:
 \pi_{n+1}(B,P)\xrightarrow{\ \cong\ }H_{n+1}(B,P).
\]
Nie używamy tutaj części lematu o relacjach.
Ponieważ $B$ jest ściągalny, oba długie ciągi pary dają
izomorfizmy brzegowe. Zgodność odwzorowań Hurewicza
z brzegiem, sprawdzona po ich definicji, daje diagram
\[
\begin{tikzcd}[column sep=large,row sep=large]
 \pi_{n+1}(B,P) \arrow[r,"\partial","\cong"']
  \arrow[d,"h_{n+1}^{\mathrm{rel}}" left,"\cong" right]
 &\pi_n(P) \arrow[d,"h_n"]\\
 H_{n+1}(B,P) \arrow[r,"\delta","\cong"']
 &H_n(P).
\end{tikzcd}
\]
Trzy zaznaczone strzałki są izomorfizmami, więc $h_n$
po prawej również jest izomorfizmem. W obu grupach bazę
stanowią sfery $\partial Q_j$, zorientowane jako brzegi
samodzielnych dysków $Q_j$. Po stronie homotopii doprowadzamy
ich punkty bazowe drogami w $P$ do jednego punktu na $\partial B$.
Wybór tych dróg nie wpływa na klasy, ponieważ $\pi_1(P)=0$.

Niech $c_P$ będzie sumą wszystkich $(n+1)$-sympleksów
triangulacji $P$, z orientacjami odziedziczonymi z $B$.
W jego brzegu wewnętrzne ściany znoszą się parami.
Pozostaje cykl $\omega_{\mathrm{out}}$ na zewnętrznej
sferze $\partial B$ oraz cykle $-\omega_j$ na sferach
$\partial Q_j$: orientacja brzegu otworu jest przeciwna
do orientacji brzegu wyjętego dysku. Zatem
\[
 \partial c_P=\omega_{\mathrm{out}}-\sum_j\omega_j,
 \qquad
 [\omega_{\mathrm{out}}]=\sum_j[\omega_j]
 \quad\text{w }H_n(P).
\]
Ponieważ $h_n:\pi_n(P)\to H_n(P)$ identyfikuje wskazane
bazy, ta sama równość zachodzi w $\pi_n(P)$:
klasa zewnętrznej sfery jest sumą klas sfer wewnętrznych.

Odwzorowanie $H|_P$ omija wszystkie wybrane punkty $z_\beta$.
Promieniowe deformacje komórek z usuniętymi punktami,
nieruchome na ich brzegach, przesuwają je więc do
odwzorowania $G:P\to W$. Na $\partial B$ niczego nie zmieniamy.
Na $\partial Q_j$ otrzymujemy odwzorowanie dołączające
$\varphi_{\beta(j)}$, poprzedzone homeomorfizmem sfery
o znaku $\varepsilon_j$. Znak ten daje odpowiednio klasę
$[\varphi_{\beta(j)}]$ albo jej odwrotność.
Uzasadnia to już pierwsza część lematu dla bazy będącej
punktem: identyfikuje $\pi_n(S^n)$ z $\Z$ przez klasę orientacji,
więc homeomorfizm zachowujący orientację reprezentuje $1$,
a odwracający ją reprezentuje $-1$.
Wynika stąd, że po złożeniu z $G$ równość sfer daje
sumę $\varepsilon_j[\varphi_{\beta(j)}]$ w $\pi_n(W)$.
Drogi doprowadzające są teraz obrazami dróg w $P$;
możemy zastąpić je drogami wybranymi przy definiowaniu
$[\varphi_\beta]$, ponieważ $W$ jest jednospójna.

Przejdźmy jeszcze przez $j_*:\pi_n(W)\to\pi_n(W,A)$.
Zewnętrzna sfera składa się z górnego dysku $F_1$,
dolnego dysku $F_0$ z przeciwną orientacją oraz bocznego
pasa o obrazie w $A$. Wyróżniony odcinek $\{*\}\times I$
tego pasa jest wysłany stale do $x_0$.
W kostkowym dodawaniu klas względnych taki pas doprowadza
oba dyski do tej samej bazy; jego przesuwanie w $A$ jest
dozwoloną homotopią względną. Otrzymujemy więc klasę
$[F_1]-[F_0]$, podczas gdy każda sfera wewnętrzna daje
$\varepsilon_jj_*[\varphi_{\beta(j)}]$. Stąd dokładnie
\[
 [F_1]-[F_0]
 =\sum_j\varepsilon_j\,j_*[\varphi_{\beta(j)}]
 \quad\text{w }\pi_n(W,A).
\]
Każda relacja po dołączeniu $(n+1)$-komórek jest zatem
całkowitą kombinacją klas ich odwzorowań dołączających.

W homologii to samo
stwierdza wzór
na brzeg komórkowy
\eqref{eq:formula-komorkowa}:
stopnie odwzorowań dołączających
są dokładnie
współczynnikami
tych relacji.

Wreszcie dla $k\geq n+2$
homotopię dysku $D^n$
można przybliżyć
komórkowo, zachowując
brzeg w $A$: jej obraz
nie musi wchodzić
do nowych $k$-komórek.
To samo dotyczy
$(n+1)$-wymiarowej
homotopii, bo $n+1<k$.
Zatem grupy $\pi_n$
nie zmieniają się.
Również względny
kompleks komórkowy
nie zmienia swoich
grup i brzegów
w stopniach $n$ i $n+1$.
\end{proof}

\begin{wniosek}[Zabijanie wybranych klas]
\label{cor:zabijanie-klas-pi}
Niech $Y$ będzie jednospójnym CW i $k\geq2$.
Dołączenie $(k+1)$-komórek wzdłuż klas $a_\beta\in\pi_k(Y)$
zmienia $\pi_k(Y)$ na iloraz przez podgrupę generowaną przez $a_\beta$,
a nie zmienia $\pi_i$ dla $i<k$.
\end{wniosek}
\begin{proof}
Oznaczmy nową przestrzeń przez $V$. Poprzedni lemat z bazą $A=Y$
i pierwszym nowym wymiarem $k+1$ mówi, że
$\pi_{k+1}(V,Y)$ jest wolna na dołączonych komórkach.
Odwzorowanie brzegowe do $\pi_k(Y)$ posyła ich generatory na $a_\beta$.
Niższe grupy względne znikają przez lokalny krok przybliżenia
komórkowego, ponieważ dyski tych wymiarów można przesunąć do $Y$.
Długi ciąg pary daje dokładnie podany iloraz i izomorfizmy w niższych stopniach.
\end{proof}

\begin{uwaga}[O przybliżeniu w lemacie]
Warunek $\pi_1(A)=0$ ma
konkretną rolę: pozwala
jednoznacznie przenosić
lokalne dyski do bazy.
Gdy $A$ nie jest
jednospójna, powstają
przesunięte przez pętle
kopie generatorów.
Wersja względnego
twierdzenia Hurewicza
bez tego warunku
wymaga ilorazu przez
działanie $\pi_1(A)$;
nie utożsamiamy wtedy
bezpośrednio całej
$\pi_n(X,A)$ z
$H_n(X,A)$.
\end{uwaga}

\begin{lemat}[Model bez niższych komórek]
\label{lem:model-wysoko-spojny}
Jeśli kompleks CW $X$
jest $(n-1)$-spójny,
$n\geq2$, to jest
homotopijnie równoważny
kompleksowi $Z$ z
jedną $0$-komórką
i bez komórek
wymiarów $1,\ldots,n-1$.
Odwzorowanie $Z\to X$
można dobrać tak,
by przenosiło
wyróżniony punkt
na $x_0$.
\end{lemat}
\begin{proof}
Zaczynamy od jednej
$0$-komórki $Z^0=\{*\}$
wysłanej do $x_0$.
Ponieważ $\pi_i(X)=0$
dla $i<n$, odwzorowanie to
jest już izomorfizmem
na tych grupach.
Dołączamy do punktu
po jednej $n$-komórce
dla wybranego zbioru
generatorów $\pi_n(X)$.
Każda taka komórka
z brzegiem sklejonym
do punktu jest $S^n$;
odwzorowanie do $X$ definiujemy
na niej reprezentantem
wybranej klasy.
Uzyskujemy surjekcję
na $\pi_n(X)$.

Teraz dla każdego
generatora jądra
$\pi_n(Z^n)\to\pi_n(X)$
dołączamy $(n+1)$-komórkę
wzdłuż reprezentującego
go odwzorowania $S^n\to Z^n$.
Jej obraz w $X$
jest zerowy w $\pi_n(X)$,
więc wybieramy
wypełnienie $D^{n+1}\to X$
i przedłużamy odwzorowanie
$Z\to X$.
Wniosek~\ref{cor:zabijanie-klas-pi}, zastosowany do $Y=Z^n$,
mówi, że dokładnie te relacje zabijają wybrane jądro; w
stopniu $n$ odwzorowanie
staje się izomorfizmem.
Algebraicznie, jeśli $q_*:\pi_n(Z^n)\to\pi_n(X)$ oznacza
otrzymaną surjekcję, to po dołączeniu komórek grupa źródła
staje się $\pi_n(Z^n)/\ker q_*$. Wypełnienia wybrane w $X$
zapewniają, że nowe odwzorowanie jest dokładnie homomorfizmem
indukowanym przez $q_*$ na tym ilorazie. Pierwsze twierdzenie
o izomorfizmie grup daje jego bijektywność.
Dołączone komórki
nie zmieniają niższych
grup dzięki
Wnioskowi~\ref{cor:pi-szkielety}.

Powtarzamy procedurę
w stopniu $k=n+1,n+2,\ldots$:
komórki $k$-wymiarowe
dołączone do punktu
realizują brakujące
elementy $\pi_k(X)$,
a komórki $(k+1)$-wymiarowe
dołączone wzdłuż klas
z jądra zabijają
dokładnie te klasy.
Dołączenie sfer $S^k$ do punktu daje surjekcję na brakujące klasy
i nie narusza już ustalonych niższych grup: nowy bukiet ma retrakcję
na poprzedni etap, a względne dyski wymiaru $<k$ można przesunąć do niego.
Dokładniej, jeśli $Y$ jest bieżącym etapem, to
$Y\hookrightarrow Y\vee\bigvee_\alpha S^k$ indukuje surjekcję
na $\pi_i$ dla $i<k$, ponieważ odpowiednia grupa względna jest zerowa.
Retrakcja ściskająca nowe sfery do punktu jest lewą odwrotnością
tej inkluzji, więc indukuje także jej iniektywność. Otrzymujemy
izomorfizm w każdym wcześniej ustalonym stopniu.
Do zabijania jądra w każdym stopniu stosujemy
Wniosek~\ref{cor:zabijanie-klas-pi} z bazą równą \emph{całemu bieżącemu etapowi}.
Nie wolno tu ponownie traktować punktu jako jedynej bazy pierwszych komórek,
bo wcześniej dołączono już komórki niższych wymiarów.
Odwzorowania dołączające przybliżamy komórkowo, dzięki czemu otrzymujemy strukturę CW;
homotopia takiego odwzorowania nie zmienia jej klasy ani możliwości wypełnienia w $X$.
Wszystkie odwzorowania do $X$ przedłużają się,
bo obrazy odwzorowań
dołączających z jądra
są zerowe w $\pi_k(X)$.
Komórki późniejszego
wymiaru nie naruszają
ustalonych już
izomorfizmów w
niższych stopniach.
Po wykonaniu
wszystkich etapów otrzymujemy CW $Z$.
Odwzorowania sfer i homotopie mają zwarte dziedziny, więc z Lematu~\ref{lem:zwartosc-cw-pi}
każda klasa i każda relacja pojawiają się już w skończonym etapie.
Stabilizacja grup oznacza zatem, że odwzorowanie $Z\to X$
indukuje izomorfizmy
na wszystkich
$\pi_k$ oraz na
$\pi_0$. Aby jawnie sprawdzić tę stabilizację, ustalmy $i$.
Każdy reprezentant w $\pi_i(X)$ został uzyskany na etapie $i$
(dla $i<n$ obie grupy są zerowe). Jeśli reprezentant w $Z$
ginie w $X$, jego zwarty obraz leży w pewnym skończonym etapie
konstrukcji. Przechodzimy do etapu nie wcześniejszego niż $i$;
tam odwzorowanie na $\pi_i$ jest już iniektywne, więc reprezentant
ma nullhomotopię w tym etapie, a tym samym w $Z$.
To daje osobno surjektywność i iniektywność odwzorowania końcowego.
Twierdzenie
\ref{thm:whitehead}
zamienia tę słabą
równoważność CW
w równoważność
homotopijną.
\end{proof}

\begin{twierdzenie}[Hurewicz: pierwszy niezerowy wymiar]
\label{thm:hurewicz-wyzszy}
Jeżeli $X$ jest
$(n-1)$-spójnym
kompleksem CW,
$n\geq2$, to
\[
 \widetilde H_i(X;\Z)=0
 \quad\text{dla }i<n,
 \qquad
 h_n:\pi_n(X,x_0)
 \xrightarrow{\ \cong\ }H_n(X;\Z).
\]
Izomorfizm jest naturalny:
dla odwzorowania opartego
$f:X\to Y$ odpowiedni
kwadrat z $f_*$ na
homotopii i homologii
jest przemienny.
\end{twierdzenie}
\begin{proof}
Lemat~\ref{lem:model-wysoko-spojny}
pozwala zastąpić
$X$ homotopijnie
równoważnym $Z$
z $Z^{n-1}=\{*\}$.
Homologia i homotopia
nie zmieniają się
przy tej operacji
(Twierdzenia
\ref{thm:homotopia-singularna}
i~\ref{thm:whitehead}).
W kompleksie komórkowym
$Z$ nie ma grup
łańcuchów w
stopniach $1,\ldots,n-1$,
więc
$\widetilde H_i(Z)=0$
dla $i<n$.

Niech
$E_n=\bigoplus_{\alpha\in I_n}
\Z[e_\alpha^n]$.
Z Lematu~\ref{lem:pierwsze-komorki-pi}
dla $A=\{*\}$ mamy
$\pi_n(Z^n)\cong E_n$.
Inkluzja
$Z^n\hookrightarrow Z$
jest surjekcją
na $\pi_n$,
a jądrem są relacje
narzucone przez
$(n+1)$-komórki.
Po lewej otrzymujemy
zatem
\[
 \pi_n(Z)
 \cong
 E_n/\langle
   [\varphi_\beta]:
   \varphi_\beta:S^n\to Z^n
   \text{ dołącza }e_\beta^{n+1}
 \rangle.
\]
Po prawej
$\partial_n^{\mathrm{CW}}=0$,
bo $C_{n-1}^{\mathrm{CW}}=0$,
więc
\[
 H_n(Z)
 \cong E_n/
     \operatorname{im}
       \partial_{n+1}^{\mathrm{CW}}.
\]
Obraz $[\varphi_\beta]$
przez $h_n$ jest
klasą brzegową
$(n+1)$-komórki:
jej współczynniki
w bazie $E_n$
to te same stopnie
odwzorowań dołączających
z Twierdzenia
\ref{thm:formula-komorkowa}.
Oba ilorazy są
więc ilorazami
\emph{tej samej}
wolnej grupy przez
\emph{ten sam}
podmoduł.
Ponieważ $h_n$
na generatorach
$[e_\alpha^n]$
przypisuje im
klasy komórkowe
$[e_\alpha^n]$,
jest izomorfizmem.

Na koniec
$h_n(f_*[u])
=(f\circ u)_*[S^n]
=f_*(u_*[S^n])$,
co dowodzi
naturalności.
W szczególności dla zbudowanej równoważności $q:Z\to X$
kwadrat z $h_n$ i $q_*$ jest przemienny. Oba poziome odwzorowania
indukowane przez $q$ są izomorfizmami, więc izomorficzność
$h_n$ dla $Z$ daje ją także dla wyjściowego $X$.
\end{proof}

\begin{twierdzenie}[Hurewicz względny dla par jednospójnych]
\label{thm:hurewicz-wzgledny}
Niech $(X,A)$ będzie parą CW,
$A\ne\varnothing$ i $A$
jednospójna. Jeśli para
jest $(n-1)$-spójna,
$n\geq2$, to
\[
 H_i(X,A;\Z)=0\quad(i<n),
 \qquad
 h_n^{\mathrm{rel}}:
 \pi_n(X,A,x_0)
 \xrightarrow{\ \cong\ }
 H_n(X,A;\Z).
\]
\end{twierdzenie}
\begin{proof}
Zbudujemy najpierw
model pary bez nowych
komórek poniżej
wymiaru $n$.
Sprawdźmy najpierw spójność przestrzeni. Każda składowa $X$
spotyka $A$, a $A$ jest łukowo spójna, więc $X$ ma tylko jedną
składową łukową. Ponieważ $n\geq2$, para jest co najmniej
$1$-spójna. Dokładność w
$\pi_1(A)\to\pi_1(X)\to\pi_1(X,A)$, przy zerowym pierwszym
wyrazie i trywialnym ostatnim zbiorze, daje $\pi_1(X)=0$.
Długi ciąg~\eqref{eq:les-pi}
oraz $\pi_i(X,A)=0$
dają izomorfizmy
$\pi_i(A)\to\pi_i(X)$
dla $i<n-1$ i
surjekcję dla
$i=n-1$.
Zaczynamy od
identyczności $A\to A$
oraz inkluzji $A\to X$.
Dołączamy $n$-komórki
wzdłuż reprezentantów
jądra
$\pi_{n-1}(A)\to
\pi_{n-1}(X)$.
Każdy taki reprezentant
wypełnia się w $X$,
więc odwzorowanie do $X$
przedłuża się na
nową komórkę.
Oznaczmy otrzymany etap przez $Z_{(n)}$: zawiera on całe $A$
i właśnie dołączone $n$-komórki. Nawias w indeksie odróżnia
etap konstrukcji od absolutnego szkieletu $Z^n$; w $A$ mogą
być komórki dowolnie wysokiego wymiaru.
W długim ciągu pary $(Z_{(n)},A)$ odwzorowanie brzegowe
$\pi_n(Z_{(n)},A)\to\pi_{n-1}(A)$
posyła każdą
charakterystyczną
$n$-komórkę na
klasę jej odwzorowania
dołączającego.
Lemat~\ref{lem:pierwsze-komorki-pi},
zastosowany do
jednospójnego $A$,
pokazuje, że
te dołączenia
zabijają dokładnie
podgrupę przez
nie generowaną,
czyli wybrane jądro.
Istotny fragment ciągu, dla $n\geq3$, ma postać
\[
 \pi_n(Z_{(n)},A)\xrightarrow{\ \partial\ }
 \pi_{n-1}(A)\longrightarrow\pi_{n-1}(Z_{(n)})
 \longrightarrow\pi_{n-1}(Z_{(n)},A)=0.
\]
Ostatnia grupa znika, bo do $A$ dołączyliśmy tylko $n$-komórki:
lokalny krok przybliżenia komórkowego przesuwa do $A$
każdy dysk mniejszego wymiaru. Obraz $\partial$ jest dokładnie
wybranym jądrem, a więc $\pi_{n-1}(Z_{(n)})$ jest ilorazem
$\pi_{n-1}(A)$ przez to jądro. Ponieważ odwzorowanie
$\pi_{n-1}(A)\to\pi_{n-1}(X)$ jest surjektywne,
otrzymujemy izomorfizm $\pi_{n-1}(Z_{(n)})\to\pi_{n-1}(X)$.
W stopniach
niższych niczego
nie zmieniają: dla $i<n-1$ znikają obie sąsiednie grupy względne
$\pi_{i+1}(Z_{(n)},A)$ i $\pi_i(Z_{(n)},A)$, więc
$\pi_i(A)\to\pi_i(Z_{(n)})$ jest izomorfizmem.
Gdy $n=2$,
wyjściowe $\pi_1(A)$
jest już zerowe
i ten krok nie
wymaga komórek.
W obu przypadkach wszystkie bieżące etapy są jednospójne:
zaczynamy od jednospójnego $A$ i dołączamy komórki wymiaru
co najmniej dwa. Pętlę można przesunąć do $A$, a tam wypełnić dyskiem.

Teraz dla $k=n,n+1,\ldots$
postępujemy jak w
dowodzie Lematu
\ref{lem:model-wysoko-spojny}:
dołączamy $k$-komórki
z brzegiem wysłanym
do $x_0$, aby
otrzymać surjekcję
na $\pi_k(X)$,
a potem
$(k+1)$-komórki
wzdłuż generatorów
jądra, aby
otrzymać izomorfizm.
Nigdy nie dołączamy
komórek wymiaru
mniejszego od $n$.
W każdym kroku wniosek o zabijaniu klas stosujemy z bieżącym,
jednospójnym etapem jako bazą. Przybliżamy komórkowo odwzorowania dołączające.
Powstaje odwzorowanie par $F:(Z,A)\to(X,A)$
równe identyczności na $A$, które
indukuje izomorfizmy
wszystkich grup
homotopii przestrzeni
oraz składowych.
Twierdzenie~\ref{thm:whitehead}
zapewnia równoważność
homotopijną $Z\simeq X$.
Porównując długie ciągi
homotopii i homologii
obu par, widzimy
ponadto, że $F$
indukuje izomorfizmy
grup \emph{względnych}:
odwzorowania na wyrazach dotyczących $A$ są identycznością,
a na wyrazach dotyczących $Z,X$ izomorfizmami.
Dla stopni co najmniej $3$ stosujemy lemat pięciu do grup abelowych.
W stopniu $2$ względne grupy są ilorazami $\pi_2(Z)$ i $\pi_2(X)$
przez obrazy $\pi_2(A)$, ponieważ $\pi_1(A)=0$;
to daje izomorfizm bez stosowania lematu pięciu do zbiorów.
Stopień $1$ jest trywialny po obu stronach, bo wszystkie trzy przestrzenie
są jednospójne. Dla homologii stosujemy zwykły lemat pięciu.

W modelu $(Z,A)$
nie ma względnych
łańcuchów komórkowych
w wymiarach $<n$,
więc $H_i(Z,A)=0$
dla $i<n$.
Lemat~\ref{lem:pierwsze-komorki-pi}
oblicza zarówno
$\pi_n(Z,A)$,
jak i $H_n(Z,A)$
jako tę samą
wolną grupę na
$n$-komórkach
modulo te same
relacje od
$(n+1)$-komórek.
Na generatorach
odwzorowanie $h_n^{\mathrm{rel}}$
przypisuje dyskowi
jego klasę orientacji,
więc jest izomorfizmem.
Naturalność
$h_n^{\mathrm{rel}}$
przenosi ten wynik
przez $F$
na parę $(X,A)$.
\end{proof}

\begin{przyklad}[Dysk względem brzegu]
\label{prz:pi-dysk-brzeg}
Dla pary $(D^{n+1},S^n)$
długi ciąg homotopii zawiera fragment
$\pi_{n+1}(D^{n+1})\to\pi_{n+1}(D^{n+1},S^n)
\xrightarrow{\partial}\pi_n(S^n)\to\pi_n(D^{n+1})$.
Ściągalność dysku zeruje obie grupy skrajne.
Dokładność daje więc zerowe jądro i pełny obraz $\partial$, czyli
\[
 \partial:
 \pi_{n+1}(D^{n+1},S^n)
 \xrightarrow{\ \cong\ }
 \pi_n(S^n)
 \qquad(n\geq1).
\]
Generator po lewej jest reprezentowany przez dysk tożsamościowy,
a jego brzegiem jest
tożsamość sfery.
W homologii dzieje się
to samo na mocy
Twierdzenia~\ref{thm:homologia-sfer};
oba ciągi łączy odwzorowanie
Hurewicza.
\end{przyklad}

\begin{wniosek}[Homologiczne kryterium Whiteheada]
\label{cor:whitehead-homologiczny}
Niech $f:X\to Y$ będzie odwzorowaniem
między jednospójnymi
kompleksami CW.
Jeśli
$f_*:H_i(X;\Z)\to H_i(Y;\Z)$
jest izomorfizmem
dla wszystkich $i$,
to $f$ jest
równoważnością
homotopijną.
\end{wniosek}
\begin{proof}
Zastępujemy $f$
odwzorowaniem komórkowym
homotopijnym z nim
i tworzymy cylinder
$M_f$ z dowodu
Twierdzenia~\ref{thm:whitehead}.
Inkluzja
$i:X\hookrightarrow M_f$
ma te same odwzorowania
na homologii co $f$
po złożeniu z
retrakcyjną deformacją
$M_f\to Y$.
Zatem w długim
ciągu homologii
pary $(M_f,X)$
wszystkie grupy
$H_i(M_f,X)$
są zerowe.
Rzeczywiście, w odcinku
$H_i(X)\xrightarrow[\cong]{i_*}H_i(M_f)\to H_i(M_f,X)
\xrightarrow{\delta}H_{i-1}(X)\xrightarrow[\cong]{i_*}H_{i-1}(M_f)$
dokładność najpierw zeruje strzałkę do grupy względnej.
Następnie daje iniektywność $\delta$, której obraz leży w zerowym
jądrze ostatniej strzałki. Grupa względna musi więc być zerowa.
W stopniu zero wystarcza surjektywność $H_0(X)\to H_0(M_f)$.
Ponieważ $X$
i $M_f\simeq Y$
są jednospójne,
długi ciąg homotopii
daje
$\pi_1(M_f,X)=0$.

Załóżmy, że pewna
wyższa względna
grupa homotopii
nie znika, i weźmy
najmniejsze $n\geq2$
z $\pi_n(M_f,X)\ne0$.
Para jest wtedy
$(n-1)$-spójna,
a $X$ jednospójna.
Twierdzenie
\ref{thm:hurewicz-wzgledny}
dałoby
$H_n(M_f,X)\cong
\pi_n(M_f,X)\ne0$,
sprzeczność.
Wszystkie względne
$\pi_n$ są więc
zerowe, a długi
ciąg homotopii
pokazuje, że
$i$, a zatem $f$,
indukuje izomorfizmy
wszystkich $\pi_n$.
Twierdzenie
\ref{thm:whitehead}
kończy dowód.
\end{proof}

\section{Obliczenia i ograniczenia twierdzeń}
\label{sec:pi-przyklady}

\begin{przyklad}[Pierwsza grupa homotopii sfery]
\label{prz:pi-sfera-hurewicz}
Z Wniosku~\ref{cor:pi-szkielety}
wiemy, że $S^n$
jest $(n-1)$-spójna.
Twierdzenia
\ref{thm:hurewicz-wyzszy}
i~\ref{thm:homologia-sfer}
dają dla $n\geq2$
\[
 \pi_i(S^n)=0\quad(1\leq i<n),
 \qquad
 \pi_n(S^n)\cong\Z.
\]
Tożsamość $S^n\to S^n$
odpowiada generatorowi,
a odwzorowanie $g:S^n\to S^n$
odpowiada liczbie
$\deg g$ z
Definicji~\ref{def:stopien-sfer}.
Zatem oparte $g$ jest homotopijne do tożsamości lub do ustalonego
odbicia sfery zachowującego punkt bazowy odpowiednio wtedy,
gdy stopień wynosi $+1$ lub $-1$.
Odbicie ma stopień $-1$, bo odwraca orientację dysku i jego brzegu.
Dla nieopartego $g$ najpierw przesuwamy obraz wyróżnionego punktu do bazy
i przedłużamy tę homotopię na sferę przez własność HEP.
Jednospójność sfery sprawia, że wybór drogi nie wpływa na wynik. W szczególności
odwzorowanie sfery stopnia
$\pm1$ jest
równoważnością
homotopijną:
wywołuje izomorfizm
na $H_n$ i $H_0$,
a inne grupy homologii
sfery są zerowe,
więc można użyć
Wniosku~\ref{cor:whitehead-homologiczny}.
\end{przyklad}

\begin{przyklad}[Kompleks Moore'a]
\label{prz:moore-hurewicz}
Niech $n\geq2$ i $m\geq2$.
Do $S^n$ dołączmy
jedną $(n+1)$-komórkę
odwzorowaniem stopnia $m$.
Powstała przestrzeń
$M(\Z/m\Z,n)$ ma
kompleks komórkowy
\[
 0\longrightarrow\Z
 \xrightarrow{\;\times m\;}\Z
 \longrightarrow0
\]
w dodatnich stopniach $n+1$ i $n$; ponadto $C_0=\Z$,
a pozostałe grupy łańcuchów są zerowe. Zatem
$\partial_n=0$, a $\partial_{n+1}$ jest mnożeniem przez $m$.
Stąd $\ker\partial_n=\Z$ i $\operatorname{im}\partial_{n+1}=m\Z$,
co daje $H_n(M)=\Z/m\Z$.
Równanie $m a=0$ dla $a\in\Z$ wymusza $a=0$,
więc $\ker\partial_{n+1}=0$ i $H_{n+1}(M)=0$.
Sprawdźmy też hipotezę spójności homotopijnej.
Model ma jeden wierzchołek i komórki jedynie w wymiarach $n$ i $n+1$.
Każde oparte odwzorowanie $S^i\to M$ dla $1\leq i<n$
przesuwamy komórkowo do punktowego $i$-szkieletu, więc
$\pi_i(M)=0$. Przestrzeń jest łukowo spójna, zatem rzeczywiście
jest $(n-1)$-spójna i Hurewicz daje $\pi_n(M)=\Z/m\Z$.
Intuicja jest
dosłowna: dołączony
dysk czyni
$m$-krotne
owinięcie sfery
zerem w pierwszej
niezerowej grupie.
\end{przyklad}

\begin{przyklad}[Dlaczego potrzeba jednospójności]
\label{prz:rp2-pi2}
Nakrycie antypodalne
$p:S^2\to\mathbb{RP}^2$
indukuje izomorfizm
$\pi_2(S^2)\cong
\pi_2(\mathbb{RP}^2)$
na mocy
Stwierdzenia~\ref{prop:pi-n-nakrycie}.
Zatem
$\pi_2(\mathbb{RP}^2)\cong\Z$.
Tymczasem
$H_2(\mathbb{RP}^2)=0$
z Przykładu~\ref{prz:homologia-rp2},
więc tutaj
$h_2:\Z\to0$
nie jest izomorfizmem.
Przestrzeń rzutowa
nie jest jednospójna:
$\pi_1(\mathbb{RP}^2)=\Z/2\Z$.
To konkretna granica
Twierdzenia~\ref{thm:hurewicz-wyzszy}.
\end{przyklad}

\begin{przyklad}[Bukiet $S^1\vee S^2$]
\label{prz:pi-bukiet-pokrycie}
Przestrzeń
$X=S^1\vee S^2$
ma $\pi_1(X)\cong\Z$.
Jej nakrycie uniwersalne
składa się z prostej
z kopią $S^2$ w każdym
punkcie całkowitym.
Generator $T$ przesunięć przenosi punkt $t\in\R$ na $t+1$
i sferę $S_k^2$ na $S_{k+1}^2$. Jego potęga $T^r$, $r\in\Z$,
przesuwa odpowiednio $t$ na $t+r$ i $S_k^2$ na $S_{k+r}^2$.
Iloraz przez te przesunięcia jest właśnie $S^1\vee S^2$.
Nad małym łukiem okręgu z dala od punktu sklejenia mamy
rozłączne przedziały prostej. Nad otoczeniem punktu sklejenia,
złożonym z krótkiego łuku i małego otoczenia bazy w sferze,
mamy rozłączne kopie tego samego otoczenia przy kolejnych liczbach
całkowitych. Tak sprawdzamy lokalnie warunek nakrycia.
Jego przestrzeń jest jednospójna: każda pętla ma obraz
w skończonym odcinku z dołączonymi skończenie wieloma sferami,
a taka przestrzeń jest jednospójna z van Kampena.
Można obliczyć homologię bez zwijania całej prostej:
mamy wierzchołki w punktach całkowitych, krawędzie między nimi
i po jednej $2$-komórce każdej sfery, dołączonej stale do jej wierzchołka.
Zatem $C_2=\bigoplus_{k\in\Z}\Z e_k^2$.
Każdy okrąg dołączający jest stały, więc wszystkie współczynniki
$\partial_2$ są zerowe; nie ma też $3$-komórek, czyli
$\operatorname{im}\partial_3=0$.
Z definicji homologii komórkowej
\[
 H_2(\widetilde X)
 \cong\bigoplus_{k\in\Z}\Z.
\]
Nakrycie nie zmienia
$\pi_2$, a
$\widetilde X$
jest jednospójna,
więc Hurewicz daje
\[
 \pi_2(X)\cong
 \pi_2(\widetilde X)
 \cong\bigoplus_{k\in\Z}\Z,
 \qquad
 H_2(X)\cong\Z.
\]
Każda podniesiona
sfera $S_k^2$
rzutuje się
na tę samą sferę
wewnątrz $X$.
Opiszmy jednak jej generator przy jednej ustalonej bazie $0$
w przestrzeni nakrycia. Najpierw wybieramy drogę $u_k$ w prostej
od $0$ do $k$. Inkluzję $S_k^2$ opartą w $k$ przenosimy
wzdłuż tej drogi do bazy $0$, używając pasa zmiany bazy
z Sekcji~\ref{sec:grupy-pi-n}. Otrzymaną klasę oznaczmy przez $b_k$.
Jej obraz przez $h_2$ jest klasą orientacji $k$-tej sfery:
po zapomnieniu bazy wprowadzenie pasa daje zwykłą homotopię
odwzorowań sfery, a homotopijne odwzorowania wyznaczają
ten sam obraz klasy $[S^2]$ w homologii.
Po zrzutowaniu droga $u_k$ obiega okrąg $k$ razy, a sama sfera
nadal przechodzi na jedyną sferę w $X$.
Naturalność Hurewicza dla nakrycia $p$ daje więc
$h_2(p_*b_k)=p_*h_2(b_k)=[S^2]$ dla każdego $k$.
Ponieważ $p_*$ jest izomorfizmem na $\pi_2$, w tych bazach
odwzorowanie Hurewicza sumuje współczynniki:
$\sum_k a_k\,p_*b_k\mapsto(\sum_k a_k)[S^2]$.
Wszystkie te sumy są skończone, zgodnie z definicją sumy prostej.
Homologia traci
informację, którą
z podniesionych
sfer reprezentuje
dana klasa.
\end{przyklad}

\begin{figure}[!htbp]
\centering
\begin{tikzpicture}[>=Stealth,line cap=round,line join=round]
 \draw[manifoldgold!80!black,line width=2pt]
  (-4.5,0)--(4.5,0);
 \foreach \x/\lab in {-3/{-1},0/{0},3/{1}}{
  \shade[inner color=white,outer color=manifoldblue!37]
   (\x,.95) circle(.95);
  \draw[manifoldblue!75!black,thick]
   (\x,.95) circle(.95);
  \draw[manifoldblue!45,densely dashed,thin]
   ({\x+.95},.95)
   arc[start angle=0,end angle=180,x radius=.95,y radius=.23];
  \draw[manifoldblue!70!black,thin]
   ({\x-.95},.95)
   arc[start angle=180,end angle=360,x radius=.95,y radius=.23];
  \draw[manifoldblue!32,densely dashed,thin]
   (\x,1.9)..controls({\x-.28},1.6)
     and({\x-.28},.5)..(\x,0);
  \draw[manifoldblue!62!black,thin]
   (\x,1.9)..controls({\x+.28},1.6)
     and({\x+.28},.5)..(\x,0);
  \fill[manifoldred] (\x,0) circle(2.8pt);
  \node[below=5pt,font=\small] at (\x,0)
   {$k=\lab$};
 }
 \node[manifoldgold!75!black,font=\small]
  at (4.12,-.35) {$\R$};
 \node[font=\large] at (-4.28,1.02) {$\cdots$};
 \node[font=\large] at (4.28,1.02) {$\cdots$};
\end{tikzpicture}
\caption{Fragment nakrycia uniwersalnego
$S^1\vee S^2$: kopia sfery w każdym
punkcie całkowitym nad pętlą okręgu.
Każda sfera daje inny generator
$\pi_2(X)$, ale wszystkie mają
ten sam obraz w $H_2(X)$.}
\label{fig:nakrycie-bukiet-sfer}
\end{figure}

\begin{uwaga}[Który Whitehead?]
Twierdzenie~\ref{thm:whitehead}
dotyczy słabych równoważności
homotopijnych CW.
\emph{Torsja Whiteheada}
jest innym, późniejszym
niezmiennikiem używanym
przy prostej homotopii
i ogólniejszych
twierdzeniach o
kobordyzmach; sama
nazwa nie czyni jej
częścią tego twierdzenia.
\end{uwaga}

\section*{Dalsza droga}
\addcontentsline{toc}{section}{Dalsza droga}
Następny naturalny blok to
formy różniczkowe,
orientacja, całkowanie,
twierdzenie Stokesa
i kohomologia de Rhama.
Po nich rozwiniemy
geometrię riemannowską:
tensor metryczny, koneksję,
symbole Christoffela,
transport równoległy,
geodezyjne, odwzorowanie
wykładnicze, twierdzenie
Hopfa--Rinowa, krzywiznę
i tensory krzywizny,
pola Jacobiego,
twierdzenia porównawcze
oraz twierdzenie
Gaussa--Bonneta.
\emph{Dopiero po tym bloku}
przejdziemy do teorii
Morse'a, uchwytów,
twierdzenia Smale'a
i chirurgii.
Nie wprowadzamy
tych późniejszych
teorii w obecnym rozdziale.

\par\smallskip
{\footnotesize\noindent\emph{Dalsza lektura:}
Allen Hatcher,
\emph{Algebraic Topology},
rozdział 4, §§4.1--4.2,
\url{https://pi.math.cornell.edu/~hatcher/AT/ATch4.pdf},
zawiera pełne rozwinięcie
grup względnych,
przybliżenia komórkowego,
twierdzeń Whiteheada
i Hurewicza.\par}
\clearpage
